\chapter{程序如何从源码变成正在运行的进程}

\section{问题与目标}

刚学完 C/C++ 时，通常已经掌握三件基本操作：

\begin{enumerate}
  \item 写一个 \code{.cpp} 文件。
  \item 用编译器生成可执行文件。
  \item 在终端运行它。
\end{enumerate}

但这些操作背后的系统机制往往仍然模糊：

\begin{itemize}
  \item 头文件到底是不是被“编译进来”。
  \item \code{g++ main.cpp -o app} 背后调用了哪些工具。
  \item \code{main} 为什么通常不是程序真正执行的第一条指令。
  \item 可执行文件里除了机器码还有什么。
  \item 操作系统如何把文件变成一个进程。
  \item 指针地址为什么每次运行可能不同。
  \item 这些机制为什么会影响性能测量。
\end{itemize}

\begin{keyidea}
程序运行并不是把源码直接交给 CPU。源码要经过预处理、编译、优化、汇编、链接、加载和运行时启动，最后才变成 CPU 可以执行的机器指令。性能工程中的许多误判，都来自没有看见这条链路。
\end{keyidea}

\section{学习本章的正确方式}

本章不是要求你背诵一串工具名称。预处理器、编译器、汇编器、链接器、加载器和运行时库当然各有细节，但第一阶段更重要的是掌握每一层的边界。所谓边界，就是这一层接收什么输入，产生什么输出，对下一层作出什么承诺，又有意隐藏了哪些信息。只记住命令而不理解边界，遇到真实问题时仍然很难定位原因。

学习这条路径时，可以把一个程序看成不断改变形态的对象。刚开始它是人写的源文本，带有注释、宏、头文件、类型名、函数名、模板和命名空间。经过预处理后，它变成一个更大的源文本，很多你没有亲手写的声明已经被展开进来。经过前端后，它不再只是字符，而是带有语法结构和类型含义的内部表示。经过优化后，很多源代码层面的形状会被打散、合并、删除或移动。经过后端后，它变成目标机器能表达的汇编动作。经过汇编器后，它变成带符号表和重定位信息的目标文件。经过链接器后，它变成操作系统能加载的可执行文件或共享库。最后，操作系统和运行时库把它放入一个进程环境中，准备栈、环境变量、动态库、线程局部存储和初始化逻辑，然后才把控制权交到用户代码附近。

这条路径上最容易出错的地方，不是术语翻译，而是把不同层的规则混为一谈。例如，C++ 说对象有类型和生命周期，但 CPU 执行的是读写寄存器和内存的指令；源代码里有一个函数名，优化后的机器码里这个函数却可能已经不存在；反汇编里显示一个地址，用户态程序看到的通常也不是物理地址；benchmark 源码里写了循环，优化器仍可能证明结果无用并把它删除。把这些层次分清，是后续学习汇编、ABI、缓存、SIMD 和性能计数器的基础。

\begin{mentalmodel}
读本章时始终问四个问题：这一层的输入是什么，输出是什么，谁负责完成转换，转换后哪些信息被保留、哪些信息被丢失。只要这四个问题清楚，后面再看工具输出就不会把现象误认为本质。
\end{mentalmodel}

\section{源码、可执行文件和进程不是同一个东西}

日常交流中常说“运行这个程序”。这句话方便，但不精确。至少有三个概念需要严格区分：源码、可执行文件和进程。

源码是程序员写下的文本。它描述的是语言层面的意图：声明变量，定义函数，创建对象，调用库，表达控制流。源码不是 CPU 的输入，甚至也不一定是编译器最终直接理解的形态，因为预处理会先展开头文件和宏。

可执行文件是存放在磁盘上的二进制文件。它包含机器码，也包含操作系统和加载器需要的元数据。它是静态对象：你不运行它时，它只是文件系统中的一段内容。它可以被复制，可以被校验，可以被反汇编，也可以被多个进程同时映射。

进程是一次运行中的实例。一个可执行文件可以被运行多次，每次运行都得到一个不同的进程。每个进程有自己的虚拟地址空间、栈、堆、文件描述符、环境变量、线程状态和内核维护的资源。两个进程可能来自同一个可执行文件，但它们的栈地址、堆对象、随机种子、文件描述符状态和调度历史都可以不同。

这个区分直接影响调试和性能分析。调试源码中的函数写法是一类问题；检查二进制里是否存在某个符号是另一类问题；解释某次运行为什么变慢，又是第三类问题。性能工程常常需要在三者之间切换：先在源码层提出假设，再在可执行文件层检查编译结果，最后在进程运行层观察时间、地址、缺页、调度和硬件事件。

\begin{warningbox}
不要说“这个源码运行很慢”就停止分析。严格说，慢的是某次构建得到的二进制，在某个操作系统、某组动态库、某种输入、某个 CPU 频率和调度环境下的一次或多次运行。把上下文说清楚，不是形式主义，而是避免性能结论失真的基本要求。
\end{warningbox}

\section{阶段边界：每一步都在改变问题}

从源码到进程不是一条单纯的“翻译”链路。每一步都在改变问题的表达方式。预处理阶段关心文本替换和条件选择；前端关心语法和语义；优化器关心等价变换；后端关心目标机器；汇编器关心指令编码和重定位记录；链接器关心符号解析和地址安排；加载器关心进程地址空间；运行时关心语言环境；CPU 关心指令语义和机器状态。把这几件事都叫“编译”会掩盖很多关键细节。

预处理器不理解完整的 C++ 语义。宏替换只是文本层面的机制，它不知道一个标识符是不是局部变量，也不知道一个表达式是否会造成未定义行为。宏因此很强，也很危险。一个宏可以让源码看起来像函数调用，却在展开后重复求值实参；一个条件编译分支可以让同一份源码在不同机器上变成不同程序。

编译器前端的任务不是提高性能，而是先确定这段文本在语言规则下是什么意思。名字查找、重载解析、模板实例化、类型检查、访问控制、常量表达式判断、对象生命周期和异常规范，都是语义问题。只有当语义确定后，优化器才知道哪些变换是允许的。性能学习不能绕过语言语义，因为优化器的自由度正是从语言规则里来的。

优化器不是“尽量保留源码形状”的工具。它的目标是在保持可观察行为的前提下生成更好的程序。可观察行为在 C++ 中有专门含义，涉及 volatile 访问、输入输出、原子操作、函数调用副作用、异常、对象构造析构等。普通局部变量的中间值，如果最终不影响可观察行为，就可能被完全删除。因此，用没有结果消费的循环做性能测试，本质上是在测试优化器是否能看穿你的代码，而不是测试算法有多快。

后端把中间表示映射到具体目标。不同目标机器有不同寄存器数量、指令格式、寻址模式、调用约定和指令代价。即使源码和中间表示相同，面向 x86-64、AArch64 或 RISC-V 的后端也会做出不同选择。以后讨论性能时，必须把“语言层算法复杂度”和“目标机器上的实现代价”分开。

汇编器和链接器处理的是二进制组织问题。汇编文本中的标签、外部符号、局部跳转和数据段引用，很多在汇编阶段还没有最终地址。目标文件里会留下重定位记录，告诉链接器哪些位置以后需要修正。链接器把多个目标文件和库合成更大的二进制，解析符号，安排段和节，生成动态链接信息。你在反汇编里看到的调用目标、地址位移和符号名字，许多都是这个阶段安排后的结果。

加载器和运行时处理的是“把文件变成运行环境”的问题。加载不是把整个可执行文件逐字节复制到内存这么简单。现代操作系统通常把可执行文件和共享库的若干段映射到进程虚拟地址空间中，按权限设置页面，准备初始栈，传入参数和环境变量。动态链接器还可能在运行前或首次调用时解析共享库符号。C++ 运行时会处理全局对象构造、异常机制、线程局部存储和退出时析构。所有这些都可能影响程序启动和第一次调用。

\begin{keyidea}
阶段越往后，源码里的名字和结构就越不可靠；阶段越往前，机器上的代价和布局就越不可见。成熟的分析方式是在合适的层问合适的问题，而不是试图用一个层的直觉解释所有现象。
\end{keyidea}

\section{为什么要从这条链路讲起}

系统课程和工程实践通常不会一上来就讲某个炫目的优化技巧，而会从程序表示、机器级程序、链接、异常控制流、虚拟内存和系统调用讲起。原因很现实：没有这些基础，后面的任何性能结论都没有解释力。

假设你写了两个函数，源码层一个看起来多做了几次判断，另一个看起来更简洁。你测量后发现前者更快。没有编译链知识时，你可能会猜测“分支不一定慢”或者“编译器很聪明”。这两个说法都太粗。你应该先看优化级别、内联情况、生成的指令、数据是否在寄存器里、循环是否被展开、分支是否被消除、函数是否被常量传播、调用是否经过 PLT、第一次运行是否触发动态链接、输入是否改变 cache 状态。也就是说，真正的解释路径必须穿过本章建立的层次。

再假设有人说：“我的指针地址每次运行不同，所以内存分配不稳定，性能不可控。”这个说法也混合了多个层。地址不同可能只是 ASLR 改变了虚拟地址布局，不等于物理内存位置一定以同样方式变化，也不等于 cache 行对齐必然改变，更不等于性能差异来自堆分配器。要分析它，需要区分虚拟地址、页、物理页、cache index、堆分配策略和测量噪声。这些都建立在“进程拥有虚拟地址空间”这个基础上。

再看第三个例子：一个 benchmark 第一次运行特别慢，后面多次运行稳定。没有本章的基础时，容易把第一次慢归因于 CPU “热身”。这可能对，也可能不对。第一次运行可能包括动态链接的延迟绑定，可能包括页面首次触碰产生的缺页，可能包括文件系统缓存，可能包括全局构造，可能包括 JIT 或插件初始化，也可能只是 CPU 频率从低功耗状态升上来。只有知道程序从文件到进程的启动过程，才能把这些可能性拆开。

因此，本章不是背景闲话，而是本册的地基。后面任何一次“看汇编”“解释地址”“写实验报告”“判断优化是否有效”，都会回到这里。

\section{从一个最小程序开始}

先看一个最小但真实的 C++ 程序：

\begin{lstlisting}[language=C++]
#include <iostream>

int main()
{
    int x = 1;
    int y = 2;
    std::cout << x + y << '\n';
    return 0;
}
\end{lstlisting}

从人的角度看，它做了三件事：

\begin{center}
\begin{minipage}{0.65\linewidth}
\begin{verbatim}
创建两个整数
计算 x + y
打印结果
\end{verbatim}
\end{minipage}
\end{center}

从系统角度看，它至少涉及：

\begin{center}
\begin{minipage}{0.82\linewidth}
\begin{verbatim}
source file
  -> preprocessor
  -> compiler frontend
  -> optimizer
  -> compiler backend
  -> assembly
  -> assembler
  -> object file
  -> linker
  -> executable ELF
  -> loader
  -> process address space
  -> runtime startup
  -> main
  -> CPU executes instructions
\end{verbatim}
\end{minipage}
\end{center}

本章的任务就是把这条路径铺开。

\section{为什么这条路径必须学}

如果最终要进入高性能计算、现代 CPU 优化或 AI CPU 算子，为什么第一个基础问题不是 AVX2，而是“程序怎么运行”？

因为很多性能现象发生在源码之外：

\begin{itemize}
  \item 编译器可能删除你以为正在测量的循环。
  \item 函数可能被 inline，二进制里不再有独立函数体。
  \item 动态链接和全局构造可能影响第一次运行时间。
  \item 不同编译参数可能生成完全不同的指令。
  \item 同一个地址表达式可能映射到栈、堆、全局区或 mmap 区域。
  \item 共享库调用、PLT/GOT、lazy binding 可能影响调用路径。
  \item 程序启动成本不是 kernel 成本。
\end{itemize}

如果你不知道这些层，就很容易把系统行为误认为算法行为。

\section{翻译单元：编译器真正看到的输入}

源文件和头文件是编辑器里的概念，翻译单元才是 C/C++ 编译的基本边界。一个翻译单元可以粗略理解为“一个源文件经过预处理后的结果”。预处理会处理头文件包含、宏展开、条件编译和注释删除。一个很短的源文件，只要包含了标准库头文件，预处理后就可能变成非常大的文本；编译器前端真正分析的是这个展开后的文本。

这个模型解释了一个常见误区：头文件通常不是被单独编译后再与源文件合并。头文件的内容会被包含进当前源文件，成为当前翻译单元的一部分。因此，头文件中的声明、内联函数、模板定义、宏和条件分支，都会影响当前源文件的编译结果。

声明和定义的区别在这里非常关键。声明告诉编译器“有这样一个名字，它的类型是什么，可以怎样使用”；定义则真正提供实体本身。函数声明可以在多个翻译单元里重复出现，只要一致即可；普通外部函数定义通常只能在整个程序里有一个；类定义可以通过头文件进入多个翻译单元，但必须保持一致；模板定义需要在实例化点可见，否则编译器无法为具体类型生成代码。很多链接错误，本质上不是链接器神秘，而是声明、定义和翻译单元边界没有处理好。

翻译单元还会影响优化视野。编译器在编译一个源文件时，通常只能直接看到当前翻译单元里的函数体和类型信息。假如一个小函数的定义在头文件里，编译器可以在调用点看到它的实现，于是更容易内联、常量传播和消除分支。假如函数只在另一个目标文件里定义，当前翻译单元只看见声明，那么普通编译流程下编译器不能随意假设这个函数的内部行为。链接时优化可以扩大这个视野，但那是额外机制，不是默认一定发生。

\begin{deepdive}
翻译单元是理解 C++ 工程性能的第一道门。你看到的“函数调用开销”，可能来自函数体不可见；你看到的“模板编译慢”，可能来自模板定义在许多翻译单元中重复实例化；你看到的“改一个头文件全项目重编译”，可能来自这个头文件被许多翻译单元包含。编译速度、二进制体积和运行性能，经常同时受翻译单元组织方式影响。
\end{deepdive}

头文件保护和一次定义规则也属于这个主题。头文件保护避免同一个头文件在一个翻译单元内被重复展开导致重复声明或重复定义。一次定义规则约束整个程序中的实体定义必须一致。初学者常把“头文件只包含一次”和“一次定义规则”混为一谈。前者解决的是一个翻译单元内部的重复包含问题；后者解决的是整个程序范围内实体定义是否一致的问题。它们相关，但不是同一件事。

当你阅读大型项目时，不要只看一个源文件。你要追问：这个函数的定义在哪里，调用点是否能看到函数体，宏是否改变了源代码形态，某个配置宏是否选择了不同实现，同一个头文件进入了多少翻译单元。后续学习高性能库时，你会看到许多实现被放入头文件，不只是为了方便发布，更是为了让编译器在调用点看到常量、类型、维度、对齐和目标指令集条件，从而生成专门化代码。

\section{固定 add\_i32：从源码到 ELF 的证据链}

本章最小贯穿函数固定为：

\begin{lstlisting}[language=C++,caption={贯穿本章的最小函数}]
extern "C" int add_i32(int a, int b)
{
    return a + b;
}
\end{lstlisting}

它足够小，小到可以手工追踪；也足够真实，因为它会经历完整的预处理、编译、汇编、目标文件、符号、重定位、链接和加载。先把每层证据列出来：

\begin{center}
\begin{tabular}{p{0.20\linewidth}p{0.30\linewidth}p{0.34\linewidth}}
\toprule
层次 & 产物 & 需要观察什么 \\
\midrule
源码 & \code{add.cpp} & C++ 函数语义、参数和返回类型 \\
预处理 & \code{add.i} & 头文件和宏展开后的翻译单元 \\
前端/IR & AST 或 IR dump & 函数、参数、返回表达式的语义结构 \\
汇编 & \code{add.s} & 参数寄存器、\code{add} 或 \code{lea}、返回寄存器 \\
目标文件 & \code{add.o} & 符号表、机器码字节、重定位记录 \\
可执行文件 & ELF executable & 段、节、入口、动态链接信息 \\
进程 & \code{/proc/<pid>/maps} & text/data/stack/heap/shared libraries 的虚拟地址 \\
\bottomrule
\end{tabular}
\end{center}

\subsection{源码语义}
\code{extern "C"} 让函数名使用 C linkage，避免 C++ name mangling。这样目标文件里更容易看到符号 \code{add\_i32}。参数 \code{a} 和 \code{b} 是两个 \code{int}，返回值也是 \code{int}。源码层还没有寄存器、栈地址或机器码；它只规定可观察语义：返回两个参数的和，若有有符号溢出则进入 C++ 未定义行为边界。

\subsection{汇编语义}
在 x86-64 System V ABI 下，前两个 \code{int} 参数通常通过 \register{edi} 和 \register{esi} 传入，\code{int} 返回值在 \register{eax}。优化后可能出现这种形状：

\begin{lstlisting}[numbers=none,caption={add\_i32 的典型优化汇编形状}]
add_i32:
  lea eax, [rdi + rsi]
  ret
\end{lstlisting}

也可能是：

\begin{lstlisting}[numbers=none,caption={另一种等价形状}]
add_i32:
  mov eax, edi
  add eax, esi
  ret
\end{lstlisting}

两者都能实现同一 C++ 语义。第一种用 \instruction{lea} 计算地址形式的整数加法，不更新 flags；第二种用 \instruction{add}，会更新 flags。若后续代码不需要 flags，编译器可以选择任意更合适的形态。这里要训练的不是背某条指令，而是知道同一源码语义可以有多个机器实现。

\subsection{目标文件：符号和重定位}
若 \code{add\_i32} 被编译成目标文件，它会有符号表项。可以用：

\begin{lstlisting}[numbers=none,caption={观察目标文件身份}]
c++ -O3 -c add.cpp -o add.o
nm -C add.o
objdump -drwC -Mintel add.o
readelf -s add.o
\end{lstlisting}

如果另一个翻译单元只声明并调用它：

\begin{lstlisting}[language=C++,caption={调用方翻译单元}]
extern "C" int add_i32(int a, int b);

int call_add()
{
    return add_i32(7, 5);
}
\end{lstlisting}

调用方目标文件在编译时还不知道 \code{add\_i32} 最终地址，只能留下一个未定义符号和重定位位置。链接器把调用点修正到实际函数地址，或通过 PLT/GOT 等动态链接机制间接跳转。看到反汇编里的 \code{call 0x0} 时，不要误以为真的调用地址 0；在目标文件里它常常只是等待链接器修正的占位。

\subsection{ELF 和 loader}
链接后的 ELF 不只是代码字节。它还包含 program headers，告诉内核和动态链接器哪些区域要映射成可执行、只读、可写；包含 dynamic section，说明依赖哪些共享库；包含入口地址，通常指向运行时启动代码，而不是直接指向 \code{main}。当 shell 调用 \code{execve} 后，内核读取 ELF header，建立新的虚拟地址空间，映射 text/data 段，准备初始栈，把 \code{argv}、\code{envp} 和 auxiliary vector 放上去，然后把控制权交给入口或动态链接器。

\begin{lstlisting}[numbers=none,caption={启动时间线的压缩版}]
shell parses command
  -> execve(path, argv, envp)
  -> kernel reads ELF header
  -> map PT_LOAD segments into virtual address space
  -> map interpreter if dynamically linked
  -> dynamic linker maps shared libraries and relocates symbols
  -> runtime startup prepares C/C++ environment
  -> global constructors run
  -> runtime calls main(argc, argv, envp)
\end{lstlisting}

这解释了为什么 \code{main} 之前也可能发生错误或耗时：共享库缺失、符号解析、TLS 初始化、全局构造、首次触页、lazy binding，都不在 \code{main} 源码里，却在进程启动路径上。

\subsection{本节验收命令}
一个合格报告至少保存：

\begin{lstlisting}[numbers=none,caption={add\_i32 证据包}]
source:
  add.cpp and caller.cpp

commands:
  c++ -E add.cpp -o add.i
  c++ -S -O0 -masm=intel add.cpp -o add_O0.s
  c++ -S -O3 -masm=intel add.cpp -o add_O3.s
  c++ -c -O3 add.cpp -o add.o
  c++ -O3 add.cpp caller.cpp -o app

binary evidence:
  nm -C add.o
  objdump -drwC -Mintel add.o
  readelf -h -l -S -s app
  ldd ./app if dynamically linked
  /proc/<pid>/maps during execution
\end{lstlisting}

学完这一节，不要求你背所有 ELF 字段，但要求你能把“源码里有一个函数”追到“目标文件里有什么符号、链接后地址如何成立、进程中它位于哪个虚拟映射”。

\section{声明、定义和链接属性：名字如何跨文件成立}

程序从一个源文件走向多个目标文件时，最核心的问题是名字如何被不同阶段理解。源码层的名字给程序员阅读，编译器前端用它做语义分析；目标文件层的符号给链接器解析；运行时的地址给 CPU 执行。声明、定义和链接属性就是把这些层连接起来的规则。

\subsection{声明告诉编译器如何使用一个名字}

声明的作用是让当前翻译单元知道某个名字的类型和使用方式。例如：

\begin{lstlisting}[language=C++]
int add_i32(int a, int b);
\end{lstlisting}

这行声明告诉编译器：存在一个名为 \code{add_i32} 的函数，它接收两个 \code{int}，返回一个 \code{int}。有了这个信息，当前翻译单元就能检查调用是否类型正确，也能生成一次调用所需的目标文件引用。但声明本身通常不提供函数体。它不能让链接器最终找到代码。

这解释了一个常见现象：只写声明，编译可以通过，链接却失败。编译阶段只需要知道如何生成引用；链接阶段需要找到引用对应的定义。如果整个程序里没有提供函数体，链接器就会报告未定义引用。

\subsection{定义提供实体本身}

定义真正创建函数体或对象。例如：

\begin{lstlisting}[language=C++]
int add_i32(int a, int b)
{
    return a + b;
}
\end{lstlisting}

函数定义会让编译器生成对应代码，并在目标文件中产生一个可以被链接器解析的符号。全局变量定义会产生存储空间。类定义、模板定义、内联函数定义又有特殊规则，因为它们经常需要进入多个翻译单元，但仍必须满足一致性要求。

下面这个例子经常导致重复定义：

\begin{lstlisting}[language=C++]
// bad.hpp
int global_counter = 0;
\end{lstlisting}

如果 \filepath{bad.hpp} 被两个 \code{.cpp} 包含，每个翻译单元都会生成一个 \code{global_counter} 定义。链接器看到多个强定义，就可能报重复定义。正确写法通常是头文件放声明，某个源文件放定义：

\begin{lstlisting}[language=C++]
// good.hpp
extern int global_counter;

// good.cpp
int global_counter = 0;
\end{lstlisting}

C++17 以后也可以在合适场景使用 \code{inline} 变量，但那是另一条明确规则，不是把定义随便放头文件就安全。

\subsection{外部链接和内部链接}

链接属性决定一个名字是否能被其他翻译单元引用。外部链接的名字可以跨翻译单元解析；内部链接的名字只在当前翻译单元内部可见。

例如：

\begin{lstlisting}[language=C++]
int public_function();

namespace {
int private_helper()
{
    return 1;
}
}

static int another_private_helper()
{
    return 2;
}
\end{lstlisting}

\code{public_function} 通常具有外部链接。匿名命名空间中的 \code{private_helper} 和 \code{static} 函数通常具有内部链接。内部链接的好处是减少名字冲突，也给优化器更多自由：既然外部目标文件不能直接引用这个符号，编译器更容易删除、内联或改写它。

性能工程中，链接属性不是纯语法细节。一个函数如果对外可见，编译器和链接器可能要保留某些 ABI 边界；如果它只在当前翻译单元内部使用，优化器可以更大胆。库设计中也常通过符号可见性控制公共 ABI 面积，避免内部函数被外部用户依赖。

\subsection{C++ 名字和二进制符号名}

C++ 支持函数重载和命名空间，因此源码里的名字不能直接作为唯一二进制符号。两个函数都叫 \code{f}，只要参数类型不同，在 C++ 中可以共存。链接器却主要处理符号名。编译器需要把 C++ 名字、参数类型、命名空间、类作用域等信息编码进符号，这就是名字改编。

例如：

\begin{lstlisting}[language=C++]
int f(int);
double f(double);
namespace math {
int f(int);
}
\end{lstlisting}

源码层看起来都有 \code{f}，目标文件符号却会不同。用 \code{nm -C} 可以反改编成人能读的名字：

\begin{lstlisting}[language=bash]
nm -C build/app | rg " f\\("
\end{lstlisting}

如果要让 C++ 函数用 C 风格符号名暴露给汇编或 C 代码，可以使用：

\begin{lstlisting}[language=C++]
extern "C" int add_i32(int a, int b);
\end{lstlisting}

注意，\code{extern "C"} 主要改变语言链接和符号命名，不会自动改变参数类型的 ABI 责任，也不会让 C++ 对象、异常、引用生命周期 magically 变得适合跨语言边界。后面互操作章节会专门强调这一点。

\begin{keyidea}
声明解决“当前翻译单元能否理解这个名字”，定义解决“程序中是否真的提供这个实体”，链接属性解决“这个名字能否跨翻译单元被找到”，符号名解决“链接器看到的二进制名字是什么”。
\end{keyidea}

\section{存储期和初始化：对象什么时候进入进程}

源码里一个对象看起来只是一个变量，但从进程角度看，它还涉及存储期、初始化时机、所在映射区域和生命周期。理解这些概念，可以避免把所有内存都简单叫作“变量所在的地方”。

\subsection{自动、静态和动态存储期}

自动存储期对象通常是函数内的普通局部变量。它们随着作用域进入和退出而创建和销毁，常见实现会把它们放在栈帧中，或者在优化后直接放进寄存器，甚至被完全消除。

静态存储期对象在程序整个运行期间存在，包括全局变量、命名空间作用域变量、函数内 \code{static} 局部变量等。它们通常与可执行文件的数据段、BSS 段、只读数据段或运行时初始化表有关。它们的初始化可能发生在进入 \code{main} 前，也可能在第一次执行到某个局部静态变量时发生。

动态存储期对象通过分配器创建，例如 \code{new}、\code{malloc} 或容器内部的堆分配。它们通常位于堆或匿名映射区域，由运行时分配器管理。动态分配对象的生命周期不由作用域自动完全决定，而由所有权和释放动作决定。

\begin{center}
\begin{tabular}{p{0.22\linewidth}p{0.33\linewidth}p{0.29\linewidth}}
\toprule
存储期 & 常见例子 & 常见观察位置 \\
\midrule
自动存储期 & 函数局部变量。 & 栈、寄存器、或优化后消失。 \\
静态存储期 & 全局变量、静态局部变量。 & data、BSS、rodata、初始化表。 \\
动态存储期 & \code{new}、\code{malloc}、容器缓冲区。 & heap、匿名 mmap、分配器管理区域。 \\
\bottomrule
\end{tabular}
\end{center}

这张表是观察线索，不是语言保证。C++ 语言不要求普通局部变量一定有栈地址，也不要求优化构建保留每个变量的独立存储。调试时如果一个变量显示为 optimized out，不一定是调试器坏了，而是优化器不再维护一个源码层变量对应的稳定位置。

\subsection{零初始化、常量初始化和动态初始化}

静态存储期对象的初始化有多个阶段。零初始化会把对象先置为零值；常量初始化可以在编译或加载时完成；动态初始化需要运行代码，可能在 \code{main} 前发生。

例如：

\begin{lstlisting}[language=C++]
int global_zero;
int global_value = 42;

int make_value()
{
    return 7;
}

int global_dynamic = make_value();
\end{lstlisting}

\code{global_zero} 可能进入 BSS 相关区域，不需要在文件里保存一份全零数据。\code{global_value} 可以放入已初始化数据。\code{global_dynamic} 需要运行函数得到初始值，启动时就可能执行额外代码。

这对性能测量很重要。如果一个 benchmark 程序有复杂全局对象，第一次运行可能包含全局初始化成本。若你用整个进程时间衡量一个小 kernel，就会把这些初始化成本混进去。更稳的做法是在 \code{main} 中显式构造 workload，把 setup 和 timed region 分开。

\subsection{局部静态变量的第一次初始化}

函数内静态局部变量常用于懒初始化：

\begin{lstlisting}[language=C++]
std::vector<int>& table()
{
    static std::vector<int> values = make_values();
    return values;
}
\end{lstlisting}

第一次调用 \code{table} 时，\code{values} 需要初始化；后续调用通常只返回已有对象。C++ 还要求局部静态初始化具有线程安全语义，这可能引入 guard 检查。对普通业务代码，这很方便；对微基准，这可能让第一次样本明显更慢。

如果你测量某个函数并发现第一次调用慢很多，要检查是否有局部静态初始化、全局对象、动态链接首次解析、内存分配器初始化或首次触页。不要简单把第一次慢解释成“CPU 还没热”。

\begin{warningbox}
对象什么时候初始化，是程序启动和第一次调用成本的一部分。性能报告如果没有区分初始化成本和稳态 kernel 成本，就很容易把运行时机制误判为算法成本。
\end{warningbox}

\section{构建产物身份：你测的是哪一个程序}

性能实验和调试实验都依赖一个朴素事实：你必须知道自己运行的是哪一个二进制。大型项目中，经常同时存在多个构建目录、多个配置、多个编译器、多个同名可执行文件和多个动态库版本。只写“运行 app”远远不够。

\subsection{二进制身份由多个因素共同决定}

一个可执行文件的身份至少包括：

\begin{itemize}
  \item 源码版本和未提交修改。
  \item 编译器名称和版本。
  \item 编译参数和宏定义。
  \item 目标架构和 ISA 选项。
  \item 链接命令、库版本和链接方式。
  \item 是否启用 LTO、PIE、调试信息、strip。
  \item 构建时间、路径和输出文件名。
\end{itemize}

这些因素任何一个变化，都可能改变机器码或运行时行为。比如 \code{-DNDEBUG} 会移除 \code{assert}；\code{-march=native} 可能启用更宽的向量指令；LTO 可能跨文件内联；不同标准库版本可能改变容器实现；PIE 和 ASLR 可能改变地址布局。

\subsection{报告中的二进制身份证}

建议每次正式实验都保存一个二进制身份证：

\begin{lstlisting}
source:
    git commit
    dirty status
    diff stat if dirty

compiler:
    clang++ --version
    g++ --version if used

build:
    full compile and link command
    build directory
    output binary path

binary:
    file output
    size output
    sha256sum if important
    ldd output for dynamic dependencies
\end{lstlisting}

对小规模实验来说，这看起来很严肃；对真实性能工程来说，这是基本卫生。没有二进制身份，实验结果很难复现，也很难判断一次变化来自源码、编译器、链接器、库还是环境。

\subsection{旧二进制问题}

最常见的低级错误是：源码改了，构建失败了，旧二进制还在，脚本继续运行旧二进制。另一个常见错误是：构建成功了，但运行命令指向另一个目录下的同名程序。第三个常见错误是：可执行文件是新的，但动态库仍然加载旧版本。

排查时可以先做三件事：

\begin{lstlisting}[language=bash]
ls -l build-release/app
file build-release/app
ldd build-release/app
\end{lstlisting}

再看构建日志和反汇编。若你刚改了目标函数，但反汇编没有变化，要么改动没有进入构建，要么优化器把差异消除了，要么你看的不是实际运行的二进制。三种情况都需要证据区分。

\begin{keyidea}
运行路径、二进制路径、动态库路径和源码版本必须能互相对上。否则，性能数字再精确，也可能只是旧程序的数字。
\end{keyidea}

\section{预处理器：文本层面的改写}

\term{预处理器}{preprocessor} 处理 include directive、宏、条件编译和注释删除。

例子：

\begin{lstlisting}[language=C++]
#define SCALE 4

int f(int x)
{
    return x * SCALE;
}
\end{lstlisting}

预处理后大致变成：

\begin{lstlisting}[language=C++]
int f(int x)
{
    return x * 4;
}
\end{lstlisting}

命令：

\begin{lstlisting}[language=bash]
mkdir -p scratch/ch01
cat > scratch/ch01/preprocess.cpp <<'EOF'
#include <iostream>

#define SCALE 4

int f(int x)
{
    return x * SCALE;
}

int main()
{
    std::cout << f(3) << '\n';
}
EOF

g++ -E scratch/ch01/preprocess.cpp -o scratch/ch01/preprocess.i
wc -l scratch/ch01/preprocess.cpp scratch/ch01/preprocess.i
\end{lstlisting}

你通常会看到 \code{preprocess.i} 远比原文件长，因为标准库头文件展开了大量声明和模板。

\begin{warningbox}
不要以为编译器看到的就是你手写的 \code{.cpp}。宏、include 和条件编译会先改写输入。性能分析遇到不符合直觉的行为时，必要时要看预处理结果。
\end{warningbox}

\section{编译器前端：从文本到语义}

预处理结果仍然是 C++ 文本。编译器前端要把它转成内部结构。

简化流程：

\begin{center}
\begin{minipage}{0.78\linewidth}
\begin{verbatim}
characters
  -> tokens
  -> parser
  -> AST
  -> semantic analysis
  -> IR generation
\end{verbatim}
\end{minipage}
\end{center}

\subsection{Token}

代码：

\begin{lstlisting}[language=C++]
int x = a + 3;
\end{lstlisting}

会被词法分析成类似：

\begin{center}
\begin{verbatim}
int
x
=
a
+
3
;
\end{verbatim}
\end{center}

\subsection{AST}

\term{抽象语法树}{Abstract Syntax Tree} 描述语法结构。上面的声明可理解成：

\begin{center}
\begin{minipage}{0.68\linewidth}
\begin{verbatim}
declaration x
  type: int
  initializer:
    binary +
      left: a
      right: 3
\end{verbatim}
\end{minipage}
\end{center}

\subsection{语义分析}

语义分析检查：

\begin{itemize}
  \item 名字是否声明。
  \item 类型是否匹配。
  \item 函数调用参数是否正确。
  \item 重载解析选择哪个函数。
  \item 模板实例化是否合法。
  \item \code{const}、引用和生命周期规则是否满足。
\end{itemize}

这一步决定了程序的 C++ 语义。后续优化必须在这个语义边界内进行。

\section{IR 和优化器：在语义边界内重写程序}

编译器通常不会直接从 C++ 生成最终机器码，而是先生成 \term{中间表示}{intermediate representation}，简称 IR。

IR 的作用是提供一个比 C++ 简单、比机器码抽象的优化平台。

例如：

\begin{lstlisting}[language=C++]
extern "C" int add_i32(int a, int b)
{
    return a + b;
}
\end{lstlisting}

LLVM IR 可能类似：

\begin{lstlisting}
define i32 @add_i32(i32 %a, i32 %b) {
entry:
  %sum = add i32 %a, %b
  ret i32 %sum
}
\end{lstlisting}

优化器会在 IR 上做：

\begin{itemize}
  \item 删除无用代码。
  \item 常量折叠。
  \item 函数内联。
  \item 循环优化。
  \item 自动向量化。
  \item 别名分析。
\end{itemize}

\subsection{一个重要例子：无用计算会消失}

如果你写：

\begin{lstlisting}[language=C++]
int sum(int n)
{
    int s = 0;
    for (int i = 0; i < n; ++i) {
        s += i;
    }
    return s;
}

int main()
{
    sum(1000000);
}
\end{lstlisting}

如果返回值没有可观察用途，优化器可能删除整个调用。  
这就是 benchmark 必须防止 dead-code elimination 的原因。

\section{后端、汇编和编译器驱动程序}

编译器后端把优化后的 IR 降低到目标机器。

后端要解决：

\begin{itemize}
  \item 指令选择：使用哪些 x86-64 指令。
  \item 寄存器分配：值放在哪些寄存器。
  \item 指令调度：指令如何排序。
  \item 调用约定：参数和返回值放哪里。
  \item 汇编输出：生成 \code{.s} 文件。
\end{itemize}

生成汇编：

\begin{lstlisting}[language=bash]
g++ -O3 -S -masm=intel scratch/ch01/add.cpp -o scratch/ch01/add.s
\end{lstlisting}

汇编是机器码的人类可读表示。CPU 最终执行的不是汇编文本，而是汇编器编码后的字节。

\subsection{你调用的命令通常是驱动程序}

很多教程会写“用编译器编译”。这句话在入门时够用，但如果要理解构建问题和性能实验，就要更精确。你在终端里调用的 \code{g++} 或 \code{clang++} 通常是驱动程序。驱动程序读取命令行参数，决定调用哪些真正的阶段工具：预处理器、编译器前端、优化器、后端、汇编器、链接器，以及目标平台需要的启动文件和运行时库。

同一个命令根据参数不同，会停在不同阶段。只预处理时输出预处理结果；生成汇编时停在汇编文本；只编译时生成目标文件；不加停止选项时继续链接生成可执行文件。理解这些停止点的意义，比记住某一个命令更重要。因为你之后定位问题时，往往要把完整构建拆开：先看预处理后宏是否按预期展开，再看汇编是否生成预期指令，再看目标文件是否有符号，再看最终可执行文件是否链接到正确库。

驱动程序还会隐式加入很多内容。你写命令时可能只给了一个源文件，但最终链接时还会加入启动代码、运行时库、标准库、动态链接器路径和默认系统库。C++ 程序尤其如此。一个看似最小的程序，只要使用标准库输入输出，就会牵涉大量库代码、初始化逻辑和动态链接信息。初学者如果以为“我的源文件只有十行，所以最终程序也只有十行机器码”，就会在第一次看反汇编时非常困惑。

优化级别、调试信息、目标 CPU、标准版本、异常、运行时类型信息、链接方式、可见性、位置无关代码，这些参数都会改变生成结果。性能实验必须记录它们，因为它们不是注释，而是程序的一部分。一个用调试构建测出的慢，不能直接代表发布构建；一个默认目标 CPU 下的汇编，不能代表启用特定指令集后的汇编；一个静态链接程序的启动和加载行为，不能直接代表动态链接程序。

\begin{warningbox}
不要把构建命令当作实验报告里的附属信息。对性能实验来说，构建命令是实验条件。没有构建命令，别人无法判断你测的是哪一个程序，更无法复现实验。
\end{warningbox}

\section{汇编器和目标文件}

\term{汇编器}{assembler} 把汇编文本变成目标文件。

命令：

\begin{lstlisting}[language=bash]
g++ -c -O3 scratch/ch01/add.cpp -o scratch/ch01/add.o
\end{lstlisting}

目标文件已经包含机器指令，但通常还不是完整程序。可以把目标文件看成一个带注释的二进制片段：里面有一些已经编码好的指令和数据，也有一些还需要链接器以后填上的位置。这个“还需要以后填上”的事实，就是重定位存在的原因。

目标文件通常按节组织内容。代码放在文本相关的节中，只读常量放在只读数据节中，已初始化全局变量放在数据节中，未初始化或零初始化全局变量放在 BSS 相关节中，调试信息放在调试节中，符号表记录名字与位置的关系，重定位节记录哪些机器码或数据位置还依赖外部符号或最终地址。节主要服务于链接器和调试工具；后面加载进进程时，加载器更关心程序头描述的可加载段。

符号表是目标文件之间互相认识的方式。一个目标文件可以定义符号，也可以引用外部符号。定义符号意味着“我这里提供这个名字对应的代码或数据”。未定义符号意味着“我需要别人提供这个名字”。链接器的基本工作之一，就是把所有未定义引用匹配到某个定义上。如果找不到，就产生未定义引用错误；如果普通强符号出现多个不兼容定义，就可能产生重复定义错误。

C++ 的符号比 C 更复杂，因为 C++ 支持函数重载、命名空间、类成员函数、模板和不同调用形式。为了让链接器区分这些实体，编译器会把源代码层面的名字编码成更详细的二进制符号名，这通常叫名字改编。你在工具输出中看到一串看起来很难读的符号，可能只是一个带命名空间和参数类型的 C++ 函数。工具的反改编选项可以把它还原成人更容易读的形式。

重定位是目标文件的另一个核心。编译一个源文件时，编译器和汇编器常常不知道某个函数或全局变量最终会被放到哪里。于是目标文件里先留下占位信息，并记录“这里以后要填入某个符号的地址或偏移”。链接器合并目标文件、安排布局、确定符号位置后，再根据重定位记录修正这些位置。位置无关代码和动态链接会让这个过程更复杂，但基本思想仍然是：某些地址相关信息不能在单个目标文件阶段完全确定。

\begin{keyidea}
目标文件不是失败的可执行文件，而是专门为链接设计的中间产物。它的价值在于保留了机器码、符号、节和重定位记录，让多个翻译单元可以独立编译，最后再组合成程序。
\end{keyidea}

从性能角度看，目标文件阶段提供了一个很好的观察窗口。你可以在还没有链接完整程序前，查看某个函数是否生成了预期指令，是否出现了不必要的栈访问，是否保留了符号，是否产生了对外部函数的调用。这种观察比只看最终可执行文件更局部，适合定位“是这个翻译单元生成得不好，还是链接后环境改变了调用路径”。

\section{链接器：把独立编译变成完整程序}

一个真实程序通常由多个目标文件和库组成。

例子：

\begin{lstlisting}[language=C++]
// add.cpp
extern "C" int add_i32(int a, int b)
{
    return a + b;
}
\end{lstlisting}

\begin{lstlisting}[language=C++]
// main.cpp
extern "C" int add_i32(int, int);

int main()
{
    return add_i32(1, 2);
}
\end{lstlisting}

构建：

\begin{lstlisting}[language=bash]
g++ -c add.cpp -o add.o
g++ -c main.cpp -o main.o
g++ add.o main.o -o app
\end{lstlisting}

\term{链接器}{linker} 负责把 \code{main.o} 里对 \code{add_i32} 的引用解析到 \code{add.o} 中的定义。

它还会处理：

\begin{itemize}
  \item 标准库。
  \item 启动代码。
  \item 静态库。
  \item 共享库引用。
  \item relocation。
  \item 符号可见性。
\end{itemize}

链接的核心不是“把文件粘起来”，而是二进制层面的名字解析、布局安排、地址修正和库选择。链接器会把多个目标文件中的节合并或安排到最终输出中，并根据重定位记录修正机器码或数据里的地址相关字段。如果是静态链接，许多引用可以在链接时完全确定；如果是动态链接，有些引用要留给动态链接器在程序加载时或首次调用时处理。

静态库和共享库体现了两种不同组合方式。静态库本质上是一组目标文件的归档，链接时需要的部分会被取出来合入最终输出。共享库则在运行时作为独立映射进入进程，多个进程可以共享只读代码页。静态链接可能让部署更简单，但二进制更大，更新库需要重新链接；动态链接可以共享库和独立更新，但启动、符号解析和版本兼容问题更复杂。

链接还引入了一个对性能很重要的边界：编译器在单个翻译单元内做优化，链接器负责跨目标文件组合。如果没有链接时优化，普通编译器很难跨越目标文件边界内联函数或做全程序分析。启用链接时优化后，编译器可以把更多中间表示保留到链接阶段，再进行跨翻译单元优化。这可能显著改变函数边界、符号可见性、内联决策和最终性能。因此，报告性能结果时必须说明是否启用了链接时优化。

\begin{deepdive}
许多看似运行时的问题，根源其实在链接阶段：链接到了错误版本的库，库顺序导致符号没解析，调试构建混入发布构建，动态库搜索路径指向旧文件，或者链接时优化改变了函数边界。遇到这类问题，应该检查目标文件符号表和链接命令，而不是只盯着调用点源码。
\end{deepdive}

\section{ELF 可执行文件：机器码之外的二进制合同}

Linux x86-64 上的可执行文件通常是 \term{ELF}{Executable and Linkable Format}。

ELF 不是“一段纯机器码”。它是一个带元数据的二进制容器，包含：

\begin{itemize}
  \item ELF header。
  \item program headers。
  \item sections。
  \item symbol table。
  \item relocation information。
  \item dynamic linking information。
  \item debug information。
\end{itemize}

查看：

\begin{lstlisting}[language=bash]
file app
readelf -h app
readelf -l app
readelf -S app
nm -C app
objdump -drwC -Mintel app
\end{lstlisting}

现在不要求你完全读懂每一项。本册只要求你建立“可执行文件是操作系统加载器能理解的一种结构化文件”这个模型；ELF、重定位、动态链接和加载器细节会在后续系统专题中继续展开。

ELF 头描述文件的基本身份，例如目标架构、文件类型、入口地址、程序头位置和节头位置。程序头描述加载器关心的段，它告诉内核或动态链接器哪些文件范围要映射成进程中的可读、可写或可执行区域。节头主要服务于链接和分析工具，描述代码节、数据节、符号表、字符串表、重定位节、调试节等。一个已经去除节头的可执行文件仍然可能运行，因为加载器主要依赖程序头；但没有足够符号和调试信息时，人类分析会困难很多。

这就解释了为什么同一个二进制可以有“运行视角”和“分析视角”。运行视角关心可加载段和入口点；分析视角关心节、符号、源码行号和调试信息。你用工具查看节表，看到的是链接器和调试器组织信息的方式；你查看程序头，看到的是加载器如何把文件映射进进程；你查看反汇编，看到的是某些可执行字节被解释成指令；你查看符号表，看到的是名字和地址之间的关系。每个视角都是真的，但它们回答的问题不同。

可执行文件也有权限结构。代码段通常可读可执行但不可写；只读数据通常可读不可写；可写数据段可读可写但不可执行；栈和堆通常不可执行。这些权限不是装饰，而是安全和调试的重要基础。把数据页设为不可执行可以阻止一些攻击；把代码页设为不可写可以防止程序随意修改自身；违反权限的访问会触发异常并由操作系统处理。后续学习段错误时，要把它理解成进程访问了当前页表和权限不允许的地址，而不是简单地说“指针坏了”。

ELF 还保存动态链接信息。动态段会记录需要的共享库、符号查找表、重定位表、初始化和终止函数等。动态链接器根据这些信息把共享库映射进进程，并处理符号绑定。第一次调用某个外部函数时，如果使用延迟绑定，调用路径可能会先进入解析逻辑，再跳到真实函数。这就是为什么一些函数第一次调用比后续调用慢，也解释了为什么严谨 benchmark 常常要预热或把第一次观测单独处理。

\begin{warningbox}
不要把 ELF 简化成“里面装着机器码”。如果只看机器码，你会漏掉入口点、段权限、动态库依赖、重定位、符号、调试信息和异常展开信息。系统性能问题往往就藏在这些元数据和运行时解释过程里。
\end{warningbox}

\section{操作系统加载程序}

当你运行：

\begin{lstlisting}[language=bash]
./app
\end{lstlisting}

shell 会请求内核创建新进程。内核和加载器大致做：

\begin{center}
\begin{minipage}{0.82\linewidth}
\begin{verbatim}
read ELF header
create process address space
map executable segments
map shared libraries if needed
prepare stack
prepare argc / argv / envp
start dynamic linker if dynamically linked
transfer control to entry point
\end{verbatim}
\end{minipage}
\end{center}

如果是动态链接程序，动态链接器还会解析共享库和符号，然后运行 C/C++ runtime 初始化。

\begin{keyidea}
\code{main} 是 C/C++ 程序员约定的入口函数，但进程真正开始执行的位置通常是 ELF header 中的 entry point，例如 \code{_start}。
\end{keyidea}

映射是理解加载的关键词。文件中的某些范围可以被映射到进程虚拟地址空间。映射并不一定意味着所有内容立刻从磁盘读入物理内存。现代操作系统通常使用按需分页：页面第一次被访问时，如果还没有物理页或内容尚未载入，就触发缺页，由内核把所需内容准备好。这样可以减少启动时不必要的读入，也让多个进程共享同一份只读代码页。

按需分页对性能解释非常重要。一个程序第一次访问大数组时可能变慢，不一定是算法突然变差，而可能是页面首次触碰导致缺页、零页分配或物理页建立。第一次执行某段很少用的库代码也可能触发指令页载入。后续运行同一路径时，页面可能已经在内存中，文件系统缓存也可能已经热起来。若不区分这些现象，就会把加载和内存管理成本误记到算法本身。

动态链接器的工作也属于加载路径的一部分。它要找到共享库文件，检查依赖，映射库段，处理符号和重定位，运行初始化函数。库搜索路径、运行时路径、系统缓存和环境变量都可能影响它找到哪个库。一个程序在开发机能运行，在另一台机器上启动失败，常常不是源码问题，而是动态库版本、路径或 ABI 不匹配。

\section{进程和虚拟地址空间}

\term{进程}{process} 可以理解为：

\begin{center}
\begin{minipage}{0.75\linewidth}
\begin{verbatim}
running program instance
  + virtual address space
  + OS-managed resources
\end{verbatim}
\end{minipage}
\end{center}

进程拥有：

\begin{itemize}
  \item 虚拟地址空间。
  \item 文件描述符。
  \item 当前工作目录。
  \item 信号处理状态。
  \item 一个或多个线程。
  \item 权限和资源限制。
\end{itemize}

简化的虚拟地址空间：

\begin{center}
\begin{minipage}{0.5\linewidth}
\begin{verbatim}
high address
  stack
  ...
  shared libraries / mmap
  ...
  heap
  bss
  data
  rodata
  text
low address
\end{verbatim}
\end{minipage}
\end{center}

这只是模型。真实地址受 ASLR、动态链接、内核配置和运行环境影响。

\section{入口点不是 \code{main}}

查看入口点：

\begin{lstlisting}[language=bash]
readelf -h app | rg "Entry point"
objdump -d -Mintel app | sed -n '/<_start>:/,+40p'
nm -C app | rg " main$|_start"
\end{lstlisting}

你会看到 \code{_start} 一类符号。它最终会调用运行时库，再调用你的 \code{main}。

这对性能测量很重要：

\begin{itemize}
  \item 程序启动成本不是函数成本。
  \item 动态链接和首次解析可能影响第一次运行。
  \item 全局对象构造可能在 \code{main} 前发生。
  \item benchmark 应该把 setup 和 timed region 分开。
\end{itemize}

\section{运行时启动：\code{main} 之前已经发生了很多事}

C++ 程序员习惯把 \code{main} 看成入口，这是语言层面的入口，不是进程第一条用户态指令。从操作系统转入用户态后，通常先执行运行时启动代码。启动代码从初始栈中取出参数、环境变量和辅助向量，准备运行时环境，调用全局构造相关逻辑，设置退出处理，然后调用 \code{main}。当 \code{main} 返回后，运行时还要处理返回值、执行退出回调和全局析构，最后通过系统调用结束进程。

全局对象构造是 C++ 特有的重要成本来源之一。带构造函数的全局对象、静态存储期对象、某些库内部对象，都可能在进入 \code{main} 前初始化。它们可能分配内存、打开文件、注册回调、初始化互斥量或触发其他库逻辑。如果 benchmark 把整个程序运行时间当作某个小函数的成本，就会把这些启动动作全部混入结果。

线程局部存储也可能参与启动和访问成本。线程局部变量给每个线程提供独立实例，这对并发程序很有用，但它需要运行时和 ABI 共同支持。不同类型的线程局部访问模型有不同代价。初学阶段不需要掌握全部细节，但要知道：看似普通的全局或线程局部变量访问，在机器层可能不是一条简单内存读写。

异常机制和栈展开信息也在二进制和运行时中占有位置。即使正常路径没有抛出异常，支持异常的 C++ 程序也可能带有展开表和处理逻辑。编译选项是否启用异常，会影响代码生成、二进制体积和调用边界。性能分析时，如果一个项目关闭异常或运行时类型信息，它生成的二进制和另一个默认项目可能不具有可比性。

标准库初始化同样不可忽略。输入输出流、区域设置、内存分配器、同步设置等，都可能在启动或首次使用时带来成本。初学者常用标准输出打印调试信息，但打印本身涉及库、缓冲、系统调用和终端设备，代价远高于普通整数运算。做微基准测试时，计时区域内应避免输出；需要验证结果时，可以在计时外消费结果。

\begin{mentalmodel}
\code{main} 是你代码的入口，不是进程世界的起点。严谨分析要把“程序启动”“运行时初始化”“测试数据准备”“被测核心区域”“结果校验”“进程退出”分开。只有被测核心区域才应该代表算法或 kernel 的性能。
\end{mentalmodel}

\section{进程资源：运行的不只是指令}

进程不只是指令流。一个进程是操作系统管理的资源集合。它有虚拟地址空间，有一个或多个线程，有文件描述符表，有信号处理状态，有当前工作目录，有权限和资源限制，有调度属性，有环境变量。一个程序运行得快或慢，可能与这些资源状态有关。

如果只记住“程序有一段机器码”，后面所有系统问题都会混层。文件找不到可能不是 \code{open} 语义问题，而是当前工作目录错了；benchmark 波动可能不是算法问题，而是 \code{OMP\_NUM\_THREADS} 或 CPU 亲和性不同；服务重启失败可能不是业务逻辑问题，而是 fd 泄漏、权限、\code{RLIMIT\_NOFILE} 或 \code{umask}。第一册必须先把进程基础合同讲清楚：一个运行实例到底带着哪些身份和资源进入 \code{main}。

\begin{center}
\small
\begin{tabular}{p{0.22\linewidth}p{0.34\linewidth}p{0.30\linewidth}}
\toprule
进程资源 & 它回答的问题 & 常见事故 \\
\midrule
\code{argv} & 本次运行被传入哪些参数 & 脚本参数和手工命令不一致 \\
\code{envp}/环境变量 & 动态库、线程数、allocator、ISA 分发怎样被配置 & 两台机器同命令不同路径 \\
当前工作目录 & 相对路径从哪里开始解析 & 输入、配置或模型文件读错 \\
fd table & 进程持有哪些打开对象 & fd 泄漏、标准输出阻塞、close-on-exec 漏掉 \\
权限和 \code{umask} & 能否访问文件，创建文件默认权限是什么 & 本地成功，服务账号失败 \\
资源限制 & fd、栈、core、地址空间等上限是多少 & 高并发下突然 \code{EMFILE} 或无 core dump \\
父子关系和退出状态 & 谁启动了谁，结果如何回收 & 僵尸进程、脚本忽略失败 \\
\bottomrule
\end{tabular}
\end{center}

\subsection{参数、环境和工作目录：同一个二进制可以跑成不同程序}

\code{argv} 是命令行参数，不是源码常量。Shell、脚本、测试框架、IDE、systemd、容器入口和在线服务管理器都可能传入不同参数。报告中写“运行了 \code{app}”不够，必须写完整命令。若程序根据参数选择输入规模、线程数、配置文件或后端实现，少一个参数就可能变成另一个实验。

环境变量会影响程序行为。动态库搜索、线程库设置、内存分配器行为、OpenMP 线程数、BLAS/oneDNN/推理库的线程数、某些高性能库的 ISA 分发策略，都可能通过环境变量控制。两个运行命令看起来一样，如果环境变量不同，实际执行路径可能不同。高质量实验报告应记录关键环境变量，至少说明没有依赖特殊环境，或者明确列出影响性能的设置。

当前工作目录和文件路径也会影响运行。相对路径从当前工作目录解析；程序可能读取配置文件、模型文件、输入数据或动态插件。如果你在不同目录运行同一个可执行文件，加载的数据可能不同，甚至链接到的插件也可能不同。性能实验应把输入数据来源和工作目录说清楚。

\begin{lstlisting}[numbers=none,caption={参数、环境和工作目录的证据}]
cat /proc/$PID/cmdline | tr '\0' ' '
tr '\0' '\n' < /proc/$PID/environ | sort
readlink /proc/$PID/cwd
readlink /proc/$PID/exe
\end{lstlisting}

这些命令的证据边界也要清楚。\filepath{cmdline} 证明内核记录的启动参数，不证明程序内部没有再读配置文件。\filepath{environ} 证明进程启动时和运行中保留的环境，不证明所有库都使用这些变量。\filepath{cwd} 证明当前工作目录，但线程和库函数仍可能使用绝对路径、临时目录或已经打开的 fd。证据要组合起来看。

\subsection{fd table：标准输入输出也是资源，不是普通变量}

文件描述符是进程与外部世界交互的重要接口。标准输入、标准输出、标准错误通常就是前三个文件描述符：\code{0}、\code{1}、\code{2}。打开文件、socket、管道、事件对象、epoll 对象、timerfd、eventfd，都可以表现为文件描述符。输出到终端、输出到文件和输出到管道，成本和阻塞行为都可能不同。因此，把打印放进计时区域，会让测量受外部设备和内核缓冲状态影响。

\begin{lstlisting}[numbers=none,caption={fd table 的最小观察}]
ls -l /proc/$PID/fd
cat /proc/$PID/fdinfo/1
strace -f -e trace=openat,close,dup,dup2,fcntl,execve ./app
\end{lstlisting}

fd table 的关键细节有三个。第一，fd 是进程局部小整数，不是文件本身。第二，多个 fd 可能通过 \code{dup}、\code{fork} 或库封装指向同一个 open file description，共享 offset 和状态。第三，\code{execve} 默认会保留未设置 close-on-exec 的 fd。若服务在启动子进程时忘记设置 \code{O\_CLOEXEC} 或 \code{FD\_CLOEXEC}，子进程可能意外持有监听 socket、日志文件、pipe 写端或临时文件，使父进程无法正确观察 EOF、端口无法释放，或者敏感 fd 泄漏给不该拥有的程序。

\begin{center}
\small
\begin{tabular}{p{0.24\linewidth}p{0.34\linewidth}p{0.28\linewidth}}
\toprule
现象 & 可能的 fd 根因 & 需要补的证据 \\
\midrule
服务无法退出 & 子进程仍持有 pipe/socket fd & \filepath{/proc/<pid>/fd}、\code{O\_CLOEXEC} \\
日志写入卡住 & stdout/stderr 指向慢终端或阻塞 pipe & fd 链接目标、\code{strace -T} \\
\code{EMFILE} & 进程 fd 用尽 & \code{ls /proc/$PID/fd | wc -l}、\code{RLIMIT\_NOFILE} \\
读写位置错乱 & 多线程共享 open file description offset & \code{dup/fork} 证据、是否用 \code{pread} \\
\bottomrule
\end{tabular}
\end{center}

这就是为什么系统程序常把 fd 生命周期写成显式协议：谁打开，谁拥有，是否跨线程，是否跨 \code{exec}，是否允许继承，错误路径谁关闭。fd 管理不好，问题常常表现为 I/O、服务退出、测试挂住或资源泄漏，而不是“fd 管理错误”这几个字。

\subsection{权限、\code{umask} 和资源限制：本地能跑不代表服务账号能跑}

进程有真实用户 ID、有效用户 ID、组、补充组、能力和命名空间上下文。普通开发命令在自己的用户下运行，线上服务可能在专用账号、容器、chroot、只读文件系统或不同挂载命名空间中运行。一次 \code{open} 失败，可能不是路径拼错，而是权限、目录执行位、SELinux/AppArmor、只读挂载或文件系统策略不同。

\code{umask} 会改变新建文件和目录的默认权限。程序调用 \code{open(path, O\_CREAT, 0666)}，最终权限还会被 \code{umask} 去掉一些位。若可靠输出、临时文件、模型缓存或 Unix domain socket 的权限没有写清，换一个启动环境就可能失败。实验报告至少要记录当前用户、工作目录权限、目标文件权限和 \code{umask}。

资源限制同样属于基础合同。常见限制包括 \code{RLIMIT\_NOFILE}、\code{RLIMIT\_STACK}、\code{RLIMIT\_CORE}、\code{RLIMIT\_AS}、\code{RLIMIT\_NPROC}。高并发服务最容易撞到 fd 上限；深递归或大栈对象可能撞到栈限制；崩溃后没有 core dump，可能只是 \code{RLIMIT\_CORE=0}。不要把这些问题解释成 C++ 语义错误。

\begin{lstlisting}[numbers=none,caption={权限和资源限制证据}]
id
umask
ulimit -a
cat /proc/$PID/limits
namei -l /path/to/input/or/output
stat -c '%A %U %G %n' /path/to/input/or/output
\end{lstlisting}

\subsection{\code{fork}、\code{execve}、\code{waitpid} 和退出状态}

很多系统程序不是单进程孤岛。Shell、构建系统、测试框架、服务器守护进程都会启动子进程。最小模型是：\code{fork} 复制出一个子进程，子进程通常再用 \code{execve} 替换成另一个程序，父进程用 \code{waitpid} 回收退出状态。这个模型直接影响构建脚本、测试驱动、服务拉起、进程池和工具链。

\begin{lstlisting}[numbers=none,caption={fork/exec/wait 的状态机}]
ParentRunning:
  process has address space, fd table, signal dispositions, cwd and environment

Forked:
  child gets a new process id
  many resources are inherited or copy-on-write

ChildBeforeExec:
  child may adjust fd table, cwd, environment and signal mask

Execed:
  child address space is replaced by a new executable image
  selected resources such as cwd, environment and non-CLOEXEC fd remain

Exited:
  child terminates with exit code or signal

Reaped:
  parent waitpid collects status; zombie is removed
\end{lstlisting}

\code{fork} 后父子进程不是两个线程。它们有不同进程 ID 和不同地址空间；内存页面通常通过 COW 共享，直到写入才复制。fd table 里的条目会被继承，且可能引用同一个 open file description。信号处理、环境、工作目录也有继承规则。子进程 \code{execve} 后，旧地址空间消失，但当前工作目录、环境变量、部分 fd、权限和资源限制等仍然影响新程序。理解这些继承关系，才能解释“为什么子进程拿到了不该拿的 fd”“为什么脚本运行的是另一个环境”“为什么父进程看到 pipe 永远不 EOF”。

退出状态也是基础合同。C/C++ 的 \code{return 0}、\code{std::exit(0)} 和 shell 中的成功状态相连；非零退出码表示失败类别；若进程被信号终止，父进程通过 wait status 能区分退出码和终止信号。脚本忽略退出状态，会在编译失败后继续运行旧二进制；测试框架错误解释退出状态，会把崩溃当成普通失败或把失败当成成功。

\begin{lstlisting}[language=C++,caption={父进程必须解释 wait status}]
int status = 0;
pid_t child = ::waitpid(pid, &status, 0);
if (child < 0) {
    return error_from_errno(errno);
}
if (WIFEXITED(status)) {
    int code = WEXITSTATUS(status);
    // code == 0 means success by convention.
}
if (WIFSIGNALED(status)) {
    int signo = WTERMSIG(status);
    // The child died from a signal such as SIGSEGV or SIGKILL.
}
\end{lstlisting}

\subsection{进程基础报告字段}

从这一章开始，凡是涉及运行、调试或测量的报告，都应该至少能回答下面这些问题：

\begin{lstlisting}[numbers=none,caption={进程基础报告字段}]
process_identity:
  pid
  parent_pid
  executable_path
  argv
  cwd
  selected_environment
  uid_gid
  umask
  resource_limits
  open_fd_summary
  signal_dispositions_or_masks_if_relevant
  exit_status_or_terminating_signal
\end{lstlisting}

这些字段不是形式主义。它们防止最基础的错误归因：运行旧程序、读错输入、写错目录、环境变量改变库行为、fd 泄漏、权限不一致、资源限制触发、信号终止被当成普通返回。后面讲 ELF、syscall、signal、I/O、线程和 benchmark 时，都默认读者已经能建立这张进程资源底图。

调度状态也属于进程运行环境。操作系统会把线程安排到 CPU 核心上运行，也可能中断、迁移或抢占它。线程迁移会改变缓存热度；抢占会让墙钟时间变长；后台任务会争用内存带宽、共享缓存和执行资源；电源管理会改变频率。后续测量章节会系统讨论这些噪声，但本章先建立一个原则：进程不是独占机器，性能数据总是运行环境和程序共同作用的结果。

\section{虚拟地址空间：把地址理解成进程内名字}

用户态程序中的指针值通常是虚拟地址。虚拟地址可以先粗略理解为进程地址空间中的名字。CPU 执行 load 或 store 时，虚拟地址要经过地址翻译，找到对应物理位置或触发异常。操作系统和硬件共同维护这种映射。这样做带来隔离、保护、按需分页、共享映射、文件映射和内存超量承诺等能力。

虚拟地址让每个进程看起来拥有自己的连续地址空间。两个进程可以使用相同的虚拟地址而互不干扰，因为它们的页表映射到不同物理页。共享库的只读代码页又可以被多个进程映射到各自地址空间中，从而共享物理内存。写时复制让父子进程在创建后先共享页面，直到某一方写入时才复制。所有这些机制都说明：指针数值只是入口，不能直接等同于物理存储位置。

地址空间通常按用途分区。文本区域保存程序代码，只读数据保存字符串字面量和常量，可写数据保存已初始化全局变量，BSS 保存零初始化全局变量，堆用于动态分配，栈用于函数调用和局部自动对象，映射区域用于共享库、文件映射和匿名映射。这个模型是简化的，但足以帮助你解释许多现象：为什么局部变量地址靠近栈，为什么动态分配对象在堆附近，为什么字符串字面量通常不应该修改，为什么访问空指针会失败。

栈和堆的区别尤其容易被误解。栈不是“快内存”，堆也不是“慢内存”。它们都是虚拟地址空间中的区域，最终访问仍然经过缓存和内存系统。栈分配通常只是移动栈指针，管理成本低，局部性也常常好；堆分配需要分配器维护元数据，可能加锁，可能触发系统调用或页分配，管理成本更高。但一旦对象已经分配好，访问栈对象和堆对象的速度主要取决于地址是否命中缓存、访问模式是否连续、是否存在依赖和争用，而不是“栈”这个名字本身。

页面是虚拟内存管理的重要单位。常见页面大小是若干 KiB，也有大页机制。首次访问新分配的大块内存时，可能触发页面建立和清零；随机访问大数组时，可能受到 TLB 容量影响；跨页访问和缺页会显著改变延迟。第二册讲缓存、TLB 和内存层级时会深入这些问题。本章先要求你知道：地址不是孤立数字，它属于页，页属于映射，映射有权限和来源。

\begin{warningbox}
不要把“地址小”“地址大”“两次地址接近”直接解释成物理距离或性能关系。用户态看到的是虚拟地址。要讨论性能，还要考虑页面映射、对齐、cache 行、TLB、访问模式和运行时分配器。
\end{warningbox}

\section{可观察性：每个工具只能回答一类问题}

学习本章时会接触许多工具。工具越多，越需要知道它们分别回答什么问题。预处理输出回答“编译器前端看到的文本是什么”。汇编输出回答“编译器后端打算让汇编器编码哪些指令”。目标文件反汇编回答“这个目标文件中已经生成了哪些机器码和重定位”。符号工具回答“哪些符号存在，哪些未定义，哪些可见”。ELF 查看工具回答“文件格式、段、节、入口和动态链接信息是什么”。运行时映射文件回答“进程当前虚拟地址空间如何映射”。调试器回答“某次运行中寄存器、栈、内存和控制流状态如何变化”。

没有哪个工具能单独告诉你全部真相。只看源码，你看不到优化后的机器码；只看汇编，你看不到操作系统调度；只看运行时间，你看不到是否测到了错误区域；只看符号表，你看不到函数是否在热路径中执行；只看平均值，你看不到第一次运行、缺页和迁移造成的异常值。高质量分析不是迷信某个工具，而是让多个工具回答各自擅长的问题，然后把证据拼起来。

工具输出还要结合构建上下文解释。同一个源文件，用不同优化级别、不同目标 CPU、不同编译器版本、不同链接方式，会得到不同输出。不同 Linux 发行版、不同动态库版本、不同内核配置，也可能改变加载和运行行为。因此，工具输出不能脱离命令、版本和环境。你应该养成在实验目录保存命令、源码、构建日志、工具输出和测量结果的习惯。

\section{性能误判：从源码直觉跳到时间数字最危险}

本章建立的路径可以用来识别常见性能误判。

第一类误判是把未被执行的代码当作性能对象。源码里存在一段循环，不代表二进制里存在这段循环；二进制里存在指令，不代表你的输入会执行到；执行到一次，不代表它是热点。优化器删除无用代码、分支条件选择、错误的测试输入、提前返回，都可能让你测到的不是想测的路径。

第二类误判是把启动成本当作核心算法成本。一个程序运行一次需要创建进程、加载映射、初始化运行时、解析动态库、构造全局对象、准备输入数据。若被测函数只运行很短时间，这些固定成本会淹没真实差异。正确做法是把计时区域缩小到核心代码，同时保留必要的预热和结果校验。

第三类误判是把调试构建当作真实性能。调试构建为了可调试性会保留更多栈变量、减少优化、增加检查或调试信息。它适合定位功能错误，不适合代表最终性能。发布构建则可能内联、展开、向量化、删除边界检查或重排代码。两者的汇编差异可能大到像两个程序。

第四类误判是把一次测量当作结论。进程调度、频率变化、缺页、缓存状态、分支预测状态、输入数据布局、动态链接首次解析，都可能让一次结果偏离稳态。可信测量需要重复、预热、隔离、记录环境、报告分布，并解释异常值。

第五类误判是把工具名当作证据。说“我用反汇编看过”还不够，你要指出哪一个函数、哪一个构建、哪一段指令、它对应哪个源码假设。说“我用计时器测过”也不够，你要说明计时范围、重复次数、输入规模、优化级别、是否防止死代码消除、是否排除输出和初始化。证据必须能被别人复查。

\begin{keyidea}
性能结论的最低标准是：源码假设、构建产物、运行路径和测量数据能够互相印证。任何只停留在源码直觉或单次时间数字上的判断，都不够可靠。
\end{keyidea}

\section{把本章路径用于真实工程阅读}

当你拿到一个陌生 C++ 工程时，可以按本章路径建立第一轮认识。先找到构建系统，确认编译器、标准版本、优化级别、目标架构和链接方式。再找到主要可执行文件或库的入口，区分哪些是业务源码，哪些是第三方库，哪些是生成文件。然后选择一个小函数，从源码追到汇编和目标文件，确认工具链真的按你理解的方式工作。最后运行程序，观察进程映射、动态库依赖和基本资源使用。

这套流程听起来慢，但它能避免很多低级错误。你不会在调试构建上浪费一周优化；不会把旧动态库当作新代码测；不会因为函数被内联就以为符号丢失是编译失败；不会把标准输出成本计入内核性能；不会把第一次运行的加载成本当作算法波动。越是复杂项目，越要先建立这些底层事实。

阅读高性能库时，还要特别关注条件编译和 ISA 分发。许多库会根据编译参数、运行时 CPU 检测或模板参数选择不同实现。源码中你看到的那段循环，未必是当前机器真正执行的路径。可能有标量后备实现、SSE 实现、AVX2 实现、AVX-512 实现，也可能通过函数指针或间接调用在运行时选择。要判断性能，必须确认当前构建和当前 CPU 走的是哪条路径。

阅读系统库时，要关注 ABI 和二进制边界。函数参数如何传递，结构体如何返回，异常能否跨库边界传播，符号是否可见，库版本是否兼容，这些都不是纯源码问题。后面 ABI 章节会详细讲调用约定，但本章先给出位置：ABI 位于源码语言和二进制执行之间，是不同目标文件、库和语言运行时能够协作的约定。

\section{怎样解释一个构建失败}

为了把本章知识真正用起来，可以从构建失败开始练习。构建失败并不可怕，可怕的是不知道错误属于哪一层。只要能定位层次，问题通常就会变得具体。

如果错误信息说缺少分号、括号不匹配、类型不匹配、名字没有声明，问题通常在当前翻译单元的前端阶段。此时你应该检查源文件、头文件包含、命名空间、声明顺序、模板参数和宏展开。不要一上来怀疑链接器，因为链接器根本还没有机会看到一个语义正确的目标文件。

如果错误信息说某个函数或变量未定义，源文件本身又能编译成目标文件，那么问题通常进入了链接阶段。常见原因包括：只写了声明没有写定义；定义所在源文件没有加入构建；库没有出现在链接命令中；静态库顺序不正确；C 和 C++ 的符号名约定不一致；函数签名在声明和定义之间不一致；模板定义没有在实例化点可见。处理这类错误时，应该检查目标文件符号表和链接命令，而不是只盯着调用点源码。

如果程序已经链接成功，但运行时提示找不到共享库，问题通常在加载阶段。可执行文件记录了它需要哪些共享库，动态链接器要在运行环境中找到这些库。库路径、运行时搜索路径、系统缓存、容器环境和安装位置都会影响结果。源码正确、链接成功，不代表部署环境一定能加载。

如果程序能启动，但一进入某个函数就崩溃，问题可能已经到运行阶段。此时要考虑空指针、越界访问、悬垂引用、栈破坏、ABI 不匹配、库版本不兼容、未定义行为、线程竞争等。反汇编和调试器可以帮助定位崩溃指令，但崩溃根因可能在更早的源码层或链接层。

\begin{mentalmodel}
遇到错误时先问：它发生在预处理、前端、后端、汇编、链接、加载，还是运行阶段？错误所在阶段决定第一批证据应该从哪里收集。跳过阶段判断，容易把简单问题查成复杂问题。
\end{mentalmodel}

举一个典型例子：你在头文件里写了一个普通函数定义，这个头文件被两个源文件包含，最后链接时报重复定义。初学者可能以为“头文件保护失效”。但头文件保护只防止同一个翻译单元里重复包含同一个头文件，不能防止这个函数定义出现在两个不同翻译单元中。链接器看到两个目标文件都定义了同一个外部符号，于是报错。正确修法可能是把函数定义移到一个源文件中，在头文件只保留声明；或者把它改成符合规则的内联函数；或者如果是内部实现细节，给它内部链接属性。这个例子说明，同一个头文件问题，可能要用翻译单元和链接规则解释。

再举一个例子：你在 C++ 源文件里调用一个 C 库函数，头文件声明看起来正确，链接却找不到符号。原因可能是 C++ 编译器对函数名做了名字改编，而 C 库导出的符号没有这种改编。解决办法通常是在声明处使用 C 链接约定。这个问题不是函数不存在，而是二进制符号名不一致。它处在语言边界和链接边界之间，正好体现了 ABI 的重要性。

\section{怎样解释一个运行成功但结果异常的程序}

程序运行成功不代表它按你的意图运行。性能工程中更常见的问题不是程序崩溃，而是程序“看起来没问题”，但结果不可信。

第一种情况是输出结果正确，但测量对象错误。假设你想测一个小函数的成本，却把进程启动、输入构造、标准输出、内存分配和结果校验都放进计时区。程序当然可以运行成功，输出也可以正确，但时间数字混合了太多不相关成本。此时要把计时范围缩小，先准备输入和内存，再开始计时，计时结束后再验证结果。

第二种情况是二进制不是你以为的二进制。构建目录里可能同时存在旧文件和新文件，脚本可能运行了旧路径，动态库可能加载了系统版本而不是本地版本，容器里可能缺少某个编译参数。程序可以运行成功，但运行的不是你刚刚修改的实现。遇到这种情况，要检查构建时间戳、符号、反汇编、动态库依赖和运行命令。

第三种情况是优化器改变了你以为的执行路径。源码里写了一个循环，结果没有被使用，优化后循环消失；源码里写了一个函数调用，优化后被内联；源码里写了一个分支，常量传播后分支被删除。这些都可能让程序输出仍然正确，却让测量不再代表原始源码形状。性能实验必须用二进制证据确认热路径确实存在。

第四种情况是第一次运行和稳态运行混在一起。第一次调用可能触发动态绑定、页面首次触碰、内存分配器初始化、分支预测和缓存状态变化。程序长期运行时的稳态成本与第一次成本可能相差很大。报告中应明确区分冷启动、预热后、固定输入和多次重复的结果。

第五种情况是输入规模不合适。太小的输入会让函数调用、计时器开销、循环控制和启动成本占比过高；太大的输入可能引入内存带宽、缺页、换页或缓存容量效应。一个算法在小输入上的快慢，不一定代表大输入；一个 kernel 在热缓存数据上的表现，也不一定代表真实工作负载。输入规模是实验条件，不是附属说明。

\begin{warningbox}
“能运行”只是最低门槛，不是证据链的终点。高质量实验至少要证明：运行的是预期二进制，执行的是预期路径，测量的是预期区域，输入规模匹配问题，结果没有被优化器删除，环境噪声已经被记录或控制。
\end{warningbox}

\section{一个完整解释样例：从函数调用到二进制证据}

假设你写了一个函数，用来计算两个整数之和。源码层面它只有一次加法。一个初学者可能直接说：“这个函数成本就是一次整数加法。”这句话通常不完整。

首先要问函数是否真的被调用。如果调用点传入常量，优化器可能直接把结果算出来。如果返回值没有被使用，整个调用可能被删除。如果函数定义在头文件里，调用点能看到函数体，编译器很可能内联它。如果函数定义在另一个目标文件中，又没有启用链接时优化，调用可能保留下来。源码层面的“一个函数”在二进制层可能对应函数体、内联片段、常量结果，甚至什么都没有。

其次要问参数如何传递。按照常见 x86-64 System V ABI，前几个整数参数会放在寄存器中，返回值也放在寄存器中。此时一次简单加法函数可能只需要移动或相加寄存器，再返回。若函数没有内联，还要付出调用和返回的控制流成本。若调用跨共享库边界，还可能经过过程链接表。若第一次调用触发延迟绑定，第一次成本和后续成本又不同。

再次要问构建模式。调试构建可能保留帧指针、把局部变量放到栈上、减少寄存器优化，让汇编比发布构建复杂很多。发布构建可能只剩一两条指令，甚至完全没有独立函数。比较两个版本性能时，如果一个是调试构建，一个是发布构建，结论没有意义。

最后要问测量方式。如果把这个函数调用一次，然后用进程总时间判断成本，测到的主要是启动和计时器噪声。如果把函数调用放进大循环，又不消费结果，优化器可能删除循环。如果每次循环都打印结果，测到的主要是输出成本。正确测量需要让结果以编译器不能删除的方式被使用，同时把输出放在计时区域外，并重复足够多次。

这个例子说明：一行 C++ 表达式的真实性能，不只由表达式本身决定，还由编译可见性、优化级别、ABI、链接方式、加载行为和测量方法共同决定。本章的价值就在于让你知道每个问题应该落在哪一层。

\section{动态库边界：为什么同一个调用可能有不同路径}

现代 Linux 程序大量使用共享库。共享库让多个程序共享同一份库代码，减少磁盘和内存占用，也方便系统更新。但共享库同时引入了二进制边界。调用一个本地静态函数、调用同一目标文件里的函数、调用另一个目标文件里的函数、调用共享库导出的函数，在机器路径和优化机会方面都不完全一样。

如果函数在当前翻译单元可见，编译器可能内联它，也可能进行常量传播和删除无用分支。如果函数只在另一个目标文件里可见，普通编译阶段通常看不到它的实现，除非启用跨翻译单元优化。如果函数位于共享库中，编译器还要遵守动态链接和符号可替换性的规则。在一些平台和选项下，默认可见符号可能被运行时的其他定义替换，这会限制编译器对调用目标的假设。

过程链接表和全局偏移表是动态链接中常见的机制。初学阶段不需要背它们的每个字段，但要知道它们解决的问题：共享库可以被映射到不同虚拟地址，外部函数和全局对象的位置可能需要运行时确定。为了让代码在不同加载地址仍然工作，二进制会使用间接表和相对寻址。于是一次外部函数调用可能不是直接跳到最终函数，而是经过一层跳转或解析入口。

动态库边界还影响部署和复现。你编译时链接的库版本，运行时未必就是实际加载的库版本。系统路径、本地路径、运行时路径、环境变量和容器镜像都会影响最终选择。如果性能突然变化，除了看源码改动，还要检查动态库是否变了。高质量报告应记录关键库版本，尤其是标准库、数学库、线程库和高性能计算库。

共享库也影响符号可见性。库作者可以选择导出哪些符号，隐藏哪些内部实现。隐藏内部符号有助于减少符号冲突，改善优化机会，也让 ABI 表面更稳定。一个面向长期维护的库，不能随意把所有内部函数暴露为公共 ABI，因为一旦外部用户依赖这些符号，后续修改就会受限。

\begin{deepdive}
动态库是工程能力和性能复杂性的交换。它让程序复用、更新和部署更灵活，同时引入加载、符号解析、版本兼容、可见性和间接调用问题。以后分析库函数性能时，要先确认调用路径是否跨过动态库边界。
\end{deepdive}

\section{程序身份审计：确认你运行的是谁}

真实工程里，“我已经重新编译了”并不等于“我运行的是新程序”。程序身份审计就是把源码、构建命令、目标文件、可执行文件、动态库和运行时进程对应起来。没有这一步，调试和性能测量都可能建立在错误对象上。

一个程序身份至少包含六个层面。

第一，源码身份。它包括 Git commit、工作区 diff、生成文件版本、配置文件和编译选项来源。若工作区有未提交修改，只记录 commit hash 不够。若构建系统自动生成头文件，也要记录生成输入。

第二，构建身份。它包括编译器路径和版本、标准库、目标架构、优化级别、宏定义、include 路径、链接库和链接顺序。两个二进制可能来自同一源码，却因为 \code{-O0} 与 \code{-O3}、\code{-march=native}、\code{-fPIC}、LTO 或 sanitizer 不同而完全不同。

第三，文件身份。可执行文件有路径、mtime、大小、build-id、符号表、section 和动态依赖。运行脚本若指向旧路径，或者 \code{PATH} 中先找到另一个同名程序，你可能在测旧二进制。

第四，库身份。动态库在链接时和运行时都参与。链接命令里的 \code{-lfoo} 不保证运行时加载的就是你以为的 \filepath{libfoo.so}。运行时搜索路径、\code{LD_LIBRARY_PATH}、\code{RPATH}、\code{RUNPATH}、系统缓存和容器镜像都会影响选择。

第五，进程身份。程序启动后成为进程，有 PID、当前工作目录、环境变量、打开文件、内存映射、线程和资源限制。两个进程运行同一个可执行文件，也可能因为环境变量和工作目录不同而行为不同。

第六，实验身份。benchmark 还要记录输入、seed、重复次数、计时区域、raw samples 和报告版本。程序身份正确，只说明运行对象对；实验身份正确，才说明测量对象对。

\subsection{身份审计的最小命令组}

可以用下面命令建立最小审计：

\begin{lstlisting}[language=bash]
git rev-parse HEAD
git status --short
cmake --build build-release --verbose
realpath build-release/app
ls -l build-release/app
file build-release/app
readelf -n build-release/app | sed -n '1,80p'
ldd build-release/app
readelf -d build-release/app | rg 'NEEDED|RPATH|RUNPATH'
\end{lstlisting}

运行时再记录：

\begin{lstlisting}[language=bash]
pwd
env | sort > results/env.txt
./build-release/app --version
cat /proc/$(pidof app)/maps > results/maps.txt
\end{lstlisting}

这些命令不是每次都要完整手敲，但严肃报告应保存等价信息。尤其是性能数据异常时，身份审计能快速排除“测错程序”“加载错库”“运行旧文件”“工作目录不对”这些低级但常见的问题。

\subsection{build-id 和符号证据}

许多 ELF 文件包含 build-id。它是识别二进制的一种常用方式。调试器、core dump 和符号服务器可以用 build-id 把崩溃现场对应到正确的二进制和调试信息。若你分析 core dump，却使用了另一个构建的可执行文件，反汇编地址和源码行可能完全错位。

符号表也能帮助确认身份。若你期望某个函数独立存在，但 \code{nm} 或 \code{objdump} 中找不到，可能是被内联、被删除、隐藏为局部符号、链接到另一个库，或者构建根本没有包含该文件。符号不存在不一定是错误，但它必须能解释。反过来，符号存在也不代表热路径一定调用它；调用可能被绕过，或者当前输入不走那条路径。

\section{启动时间线：\code{main} 之前和之后}

理解程序生命周期，必须知道 \code{main} 不是进程的第一个动作。一个普通动态链接 C++ 程序启动时，大致经历下面时间线：

\begin{enumerate}
  \item shell 或父进程调用 \code{execve}。
  \item 内核读取 ELF header，建立初始虚拟地址空间。
  \item 内核映射可执行文件和动态链接器，准备初始栈。
  \item 动态链接器加载共享库、处理重定位和符号解析。
  \item C/C++ 运行时建立环境，处理初始化数组。
  \item 全局和静态对象的动态初始化执行。
  \item 控制流进入 \code{main}。
  \item \code{main} 返回或调用退出函数。
  \item 静态对象析构、atexit handler 和运行时清理执行。
  \item 进程退出，内核回收资源。
\end{enumerate}

不同平台、链接方式和运行时会有差异，但这条线足够帮助你避免两个误解。第一，进程启动成本不是你的核心算法成本。第二，全局对象初始化、动态库加载和首次符号解析可能在 \code{main} 前后引入额外行为。若 benchmark 只运行一个很短程序并测总 wall time，测到的很可能主要是启动和初始化。

\subsection{全局初始化为什么影响实验}

C++ 允许具有静态存储期的对象在 \code{main} 前动态初始化。例如某个全局 \code{std::vector}、全局 logger、全局 registry 或全局 benchmark 对象，可能在程序真正开始之前分配内存、打开文件、注册函数、读取环境变量。这些行为会影响启动时间，也可能影响 allocator 和 cache 状态。

性能实验应尽量避免让复杂全局初始化混入被测对象。更好的做法是显式 setup：在 \code{main} 或 runner 中创建输入、注册 kernel、分配输出，并把这些步骤放在计时区域外。若必须使用全局 registry，也要在报告中说明它属于 harness，而不是 kernel。

\subsection{退出路径也可能有成本}

程序结束时，静态对象析构、缓冲刷新、日志写入、临时文件清理和 runtime handler 都可能执行。如果用外部 \code{time ./app} 测一个短程序，退出成本也包含在总时间里。对于端到端工具，这可能是合理指标；对于微基准，这不是核心 kernel 成本。计时区域应该在进程内部明确开始和结束，并把输出和清理放在外面。

\section{从翻译单元到链接图}

大型 C++ 工程不是“很多文件一起编译”，而是许多翻译单元分别编译成目标文件，再由链接器组合。理解这一点有助于解释为什么某些优化能发生，某些不能发生。

一个翻译单元由源文件经过预处理后形成。编译器前端和优化器在普通编译阶段主要看到当前翻译单元。若函数定义在另一个翻译单元，当前编译器通常只知道声明，不知道函数体。于是它不能普通内联，也不能基于函数体做常量传播。链接器后续能解析符号，但传统链接器不做完整高级优化。启用 LTO 后，更多 IR 被带到链接阶段，跨翻译单元优化才成为可能。

这解释了几个现象。把函数定义放进头文件，调用点可见，可能更容易内联；把函数放进单独 \filepath{.cpp}，普通构建可能保留调用；开启 LTO 后，原本跨文件的调用可能被内联；动态库边界可能限制跨边界优化；符号可见性影响编译器是否能假设某个函数不会被替换。

\subsection{链接图不是调用图}

链接图描述目标文件和库之间的符号依赖；调用图描述运行时函数可能调用谁。两者相关但不同。一个目标文件依赖某个库符号，不代表每次运行都会调用该符号；一个间接调用可能在链接图中看不出具体目标；一个函数被内联后，运行时没有普通调用边，但链接图中可能仍有其他符号依赖。

性能分析中要避免把链接依赖当作热路径。 \code{ldd} 显示程序依赖某个库，不代表热点在该库；\code{nm} 显示某个函数存在，不代表当前输入执行它；反汇编显示某个调用存在，也不代表它在样本中占比高。热路径需要运行时 profiling 或至少输入路径分析支持。

\section{程序生命周期中的证据层级}

本章的所有材料可以整理成证据层级：

\begin{center}
\begin{tabular}{p{0.21\linewidth}p{0.35\linewidth}p{0.29\linewidth}}
\toprule
层级 & 典型证据 & 能回答的问题 \\
\midrule
源码层 & \filepath{.cpp}、头文件、宏、Git diff & 程序员写了什么语义意图 \\
构建层 & 编译命令、构建日志、compile database & 编译器按什么条件工作 \\
目标层 & \filepath{.o}、符号表、重定位 & 每个翻译单元产出了什么 \\
链接层 & 可执行文件、动态段、build-id & 最终二进制依赖什么 \\
加载层 & \code{ldd}、\filepath{/proc/pid/maps} & 运行时映射了哪些文件 \\
执行层 & GDB、trace、profiling、samples & 程序实际走了哪些路径 \\
测量层 & raw CSV、summary、环境记录 & 运行成本如何分布 \\
\bottomrule
\end{tabular}
\end{center}

做报告时，先判断你的结论需要哪一层证据。若结论是“源码使用了某个算法”，源码层足够；若结论是“这个函数没有被内联”，需要二进制层；若结论是“运行时调用了共享库版本”，需要加载或执行层；若结论是“该实现更快”，需要测量层，同时还要有正确性和二进制证据。

\begin{keyidea}
程序从源码到进程不是一条黑箱通道，而是一串可观察边界。高质量工程分析会选择与结论匹配的证据层级，而不是把源码直觉、反汇编片段和时间数字混在一起。
\end{keyidea}

\section{综合案例：从一份崩溃报告回到生命周期}

假设你收到一份崩溃报告：程序在用户机器上启动后立即 \code{SIGSEGV}，本地无法复现。一个没有生命周期模型的排查方式，是直接在源码里搜索可能空指针的位置。更稳的方式是按本章链路倒推。

第一步，确认运行的是哪个二进制。崩溃报告是否包含 build-id？用户机器上的可执行文件是否与你本地调试的文件一致？是否 stripped？调试符号是否匹配？如果 build-id 不一致，后续源码行和反汇编地址都可能错。

第二步，确认加载了哪些共享库。用户机器是否加载了旧版本 \filepath{libfoo.so}？\code{RPATH}、\code{RUNPATH}、\code{LD_LIBRARY_PATH} 是否改变？崩溃发生在 \code{main} 前，还是进入业务代码后？如果崩溃在动态初始化或共享库构造函数中，业务函数断点可能根本不会触发。

第三步，确认入口和初始化。崩溃是否发生在全局对象初始化、线程局部对象初始化、配置读取、日志系统启动、插件扫描、CPU dispatch 初始化？这些都可能在 \code{main} 之前或刚进入 \code{main} 时发生。报告中“启动后立即崩溃”不等于“第一行业务代码崩溃”。

第四步，确认环境。用户当前工作目录、配置文件路径、权限、环境变量、CPU ISA、内核版本、容器限制是否与本地不同？如果程序启动时加载配置或选择 ISA kernel，环境差异可能直接决定路径。

第五步，映射地址到模块。core dump 中的 \register{rip} 是运行时虚拟地址，需要通过进程映射找到它属于可执行文件、共享库、JIT 区域还是匿名映射。只有映射到正确模块，才能用对应二进制反汇编。不要把地址直接拿去本地二进制里查，ASLR 和不同构建会让地址失效。

第六步，回到源码和构建。确认崩溃指令后，再看它来自哪个源码路径、哪个编译选项、哪个目标文件。若崩溃在手写汇编或 ABI 边界，要检查参数寄存器、栈对齐和 callee-saved；若崩溃在内存 load，要检查对象生命周期和地址来源；若崩溃在间接调用，要检查函数指针、虚表和动态库版本。

这个案例说明，运行时问题经常跨越多个阶段。源码只是证据之一。真正的排查路径是：二进制身份、加载身份、进程环境、崩溃地址、反汇编、源码语义。跳过前面阶段，可能会在错误源码上调试正确问题。

\section{综合案例：性能比较中的“旧库”问题}

再看一个性能案例。团队说新算法比旧算法慢，但反汇编显示新算法热循环更短。最后发现运行时加载的是系统安装的旧共享库，而不是刚编译的本地库。

这个问题在生命周期链路中很典型。编译阶段生成了新的 \filepath{libfast.so}；链接阶段可执行文件记录了需要 \code{libfast.so}；加载阶段动态链接器按照搜索路径找到了另一个目录下的旧库；执行阶段自然没有运行新代码；测量阶段却把结果归因到新算法。

排查证据包括：

\begin{itemize}
  \item \code{ldd ./app} 显示实际库路径。
  \item \code{readelf -d ./app} 显示 \code{RPATH}/\code{RUNPATH}。
  \item \filepath{/proc/<pid>/maps} 显示运行时映射。
  \item \code{nm -D} 或 build-id 显示库版本。
  \item 反汇编旧库和新库，确认热循环差异。
\end{itemize}

修复可能是设置正确的运行时路径、调整安装目录、使用绝对路径测试、清理系统旧库，或在程序启动时打印实现版本。无论哪种修复，报告都应说明“之前测到的是旧库”，并废弃原性能结论。

\begin{warningbox}
性能比较必须证明运行时加载的是预期二进制。只证明源码改了、构建过了、链接过了，还不够。
\end{warningbox}

\section{深入讲解：为什么“程序”这个词容易误导}

日常语言里，我们会说“运行这个程序”“这个程序有 bug”“这个程序很慢”。但在工程分析里，“程序”这个词太粗。它可能指源码项目、某个目标文件、某个可执行文件、某个正在运行的进程、某个动态库版本、某次 benchmark 配置，甚至某个运行路径。不同含义混在一起，就会产生大量误判。

例如“这个程序很慢”可能有至少六种解释。第一，源码算法复杂度高。第二，编译器没有生成预期优化。第三，链接到了慢版本库。第四，程序启动和初始化成本高。第五，当前输入触发了冷路径。第六，benchmark 把 I/O 或分配放进计时区域。每一种解释对应不同证据和不同修复。如果不拆开“程序”，就只能凭经验猜。

再如“这个程序崩溃”也不够精确。崩溃可能发生在动态链接器加载共享库时，可能发生在全局对象初始化时，可能发生在 \code{main} 开始前，可能发生在某个线程里，可能发生在退出析构阶段。崩溃地址属于哪个模块、哪个映射、哪个 build-id，决定后续怎么找源码和反汇编。

因此，本章建议在报告中避免单独使用“程序”作为分析对象。应该写：

\begin{itemize}
  \item 源码项目：哪个 commit，是否 dirty。
  \item 构建产物：哪个路径，哪个 build-id，哪个编译参数。
  \item 运行实例：哪个 PID，哪个环境，哪些动态库映射。
  \item 执行路径：哪个函数、哪个输入、哪个分支、哪个实现版本。
  \item 测量对象：哪个 timed region，哪些样本，哪个 summary。
\end{itemize}

这不是文字洁癖，而是工程精度。对象命名越准确，证据越容易匹配；对象命名越模糊，调试越容易绕圈。

\section{深入讲解：构建系统也是程序语义的一部分}

初学者常把构建系统当成“把代码编译出来的工具”，而不是程序行为的一部分。实际上，构建系统决定哪些源文件进入翻译单元，哪些宏被定义，哪些 include 路径优先，哪些库被链接，哪些编译选项启用，哪些生成文件参与构建。这些都会改变最终程序。

一个宏可能改变结构体布局；一个 include 路径可能让你使用了另一个版本头文件；一个链接库顺序可能改变解析到的符号；一个 \code{-DNDEBUG} 可能删除断言；一个 \code{-ffast-math} 可能改变浮点语义；一个 \code{-fno-exceptions} 可能改变异常边界；一个 \code{-march=native} 可能生成不能在其他机器运行的指令。这些都不是构建系统外部的“小配置”，而是程序身份的一部分。

因此，工程报告必须保存构建命令。CMake、Make、Ninja、Bazel 等工具隐藏了很多命令细节。你要能导出 verbose build 或 compile database。只写“用 Release 构建”不够，因为不同项目的 Release flags 可能不同；同一个项目不同机器上的 toolchain file 也可能不同。

构建系统还影响增量构建正确性。某个头文件变了但目标文件没重新编译，某个生成文件没更新，某个旧目标文件残留在 build 目录，都可能造成“源码看起来对，二进制不对”。遇到诡异问题时，clean build 是一种验证，但不要把 clean build 当作长期解决方案；长期应修复依赖关系和构建规则。

\section{深入讲解：动态链接是运行时选择}

动态链接的核心复杂性在于，最终绑定不只发生在编译时。可执行文件记录需要某些共享库和符号，动态链接器在运行环境中找到库、映射库、处理重定位，并可能在第一次调用时解析函数地址。运行时环境参与了程序身份。

这对性能和调试都有影响。第一次调用某个外部函数可能触发 lazy binding，导致第一次比后续慢；不同库版本可能有不同实现；符号可见性和 interposition 可能限制编译器优化；容器和系统环境可能加载不同库。高质量 benchmark 应尽量预热或固定动态链接状态，并记录库路径。

动态链接还带来 ABI 稳定问题。共享库的公共符号一旦被外部使用，修改签名或结构体布局就可能破坏旧程序。库作者需要控制导出符号，隐藏内部实现，使用版本脚本或命名策略。调用者则需要记录库版本，避免把系统升级造成的变化误判为源码性能变化。

第一册不要求你成为动态链接专家，但要求你养成两个习惯：看到外部调用时想到 PLT/GOT 和库版本；做实验时记录实际加载的共享库。这样后续分析 ABI 和互操作时，才不会把运行时选择误认为编译时事实。

\section{地址空间观察报告应该怎样写}

本章有一个观察地址空间的实验。这个实验的目的不是让你背某次运行中打印出的具体地址，而是训练解释方式。一份合格报告至少应包含四类信息。

第一类是构建和运行条件。报告要说明编译器、优化级别、是否带调试信息、系统类型、是否启用地址随机化、运行次数和运行命令。地址空间布局和符号位置可能受这些条件影响。没有条件说明，地址比较没有意义。

第二类是对象分类。你应该把函数地址、已初始化全局变量、零初始化全局变量、字符串字面量、栈上局部变量、堆分配对象、容器内部缓冲区分别归类。不要只写“地址不同”。要说明这些对象分别属于代码、可写数据、BSS、只读数据、栈、堆或堆内部二次分配。容器对象本身如果在栈上，它的控制字段在栈上；容器管理的元素缓冲区通常在堆上。这种区分非常重要。

第三类是跨运行比较。你可以比较同一对象类别在两次运行中的地址是否变化，也可以比较同一次运行中不同类别地址的大致关系。若地址变化，应解释这可能来自地址空间布局随机化、动态库映射、堆分配状态或栈位置变化。若某些相对关系稳定，也要谨慎说明它只是当前环境下的观察，不是 C++ 语言保证。

第四类是与映射文件对应。只打印地址还不够，最好结合进程映射信息，把地址落到某个映射区间中。一个地址如果落在可执行文件的可执行映射内，可以解释为代码；落在可写映射内，可能是全局可写数据；落在栈映射内，通常是栈对象；落在匿名可写映射内，可能与堆或内存分配器有关。映射区间的权限也很重要，可读、可写、可执行权限能帮助解释为什么某些写入会崩溃。

\begin{keyidea}
地址实验的目标不是记住地址，而是把“对象类别、虚拟地址、映射区间、权限、运行差异”连起来。能建立这种连接，后续学习栈帧、堆分配、页面、TLB 和缓存时才不会断层。
\end{keyidea}

\section{从本章到后续章节的知识接口}

本章结束后，后续章节会不断复用这里的概念。第二章讲 CPU 如何执行指令时，会把可执行文件中的机器码看成指令字节，再追踪取指、译码、执行和提交。第三章讲数据、内存和指针时，会把虚拟地址空间进一步细化到对象布局、对齐、数组寻址和生命周期。工具章节会把这里提到的预处理、反汇编、符号和调试命令组织成日常工作流。

进入汇编部分时，本章的链接和 ABI 背景会变得更重要。你看到寄存器传参、栈对齐、返回地址、调用者保存和被调用者保存规则时，要记得它们不是某个编译器的随意选择，而是二进制模块互相调用所需的约定。没有这些约定，不同源文件、不同库甚至不同语言写出的函数就无法在机器码层可靠协作。

进入测量部分时，本章的加载、运行时启动和进程资源背景会变得更重要。你需要把计时区域与程序启动分离，把动态链接和首次缺页与稳态执行分离，把输出、分配、输入准备和结果验证与核心 kernel 分离。否则，测量章节再讲统计方法也救不了错误的实验边界。

进入第二册的缓存、SIMD 和 AI 算子时，本章依然有用。高性能 kernel 往往以库函数、模板实例化、ISA 分发和动态链接的形式出现。你必须确认当前机器执行的是哪一个实现，确认编译器是否生成了目标指令，确认数据是否按预期布局，确认测量没有混入加载和调度噪声。也就是说，越往高性能方向走，越不能忘记最基础的“源码如何变成进程”。

\section{程序身份审计：先确认你在分析什么}

很多排查失败不是因为缺少高级知识，而是因为一开始就分析错了对象。源码窗口里打开的是新代码，终端里运行的是旧可执行文件；当前目录下有一个库，动态链接器实际加载的是系统库；命令行参数写在笔记里，真实运行脚本又加了另一组参数；编译器输出了目标文件，最终可执行文件却链接了缓存中的旧对象。这些情况都属于程序身份不清。

程序身份审计要回答五个问题。第一，源码身份是什么，包括 Git 提交、工作区是否干净、关键生成文件是否来自当前源码。第二，构建身份是什么，包括编译器、编译参数、宏定义、include 路径、链接库路径和构建类型。第三，二进制身份是什么，包括可执行文件绝对路径、文件时间、校验和、架构、调试信息和 build id。第四，运行身份是什么，包括命令行、工作目录、环境变量、输入文件、配置文件、权限和容器环境。第五，加载身份是什么，包括共享库实际路径、RPATH/RUNPATH、\code{LD_LIBRARY_PATH}、符号解析和插件路径。

只有这五类身份一致，后面的调试和测量才有基础。若某个身份不清，结论应降级。例如你可以写“当前反汇编来自 \filepath{build/app}，但尚未确认线上进程是否使用同一 build id”，而不要写“线上执行了这段指令”。工程报告中承认身份缺口，比假装确认更可靠。

\subsection{最小身份记录}

一次可复现实验至少应保存下面信息：

\begin{lstlisting}
source:
    git commit
    git status --short

build:
    compiler version
    full compile and link commands
    build type and important macros

binary:
    absolute path
    sha256 or build-id
    file output
    ldd output if dynamically linked

run:
    working directory
    command line
    environment differences
    input and config paths
\end{lstlisting}

这些信息看似普通，却能排除大量伪问题。比如 benchmark 结果突然变化，先比对二进制校验和和动态库路径，可能立刻发现新旧版本混用。崩溃只在线上发生，先比对 build id 和符号文件，才能保证 core dump 对应正确源码。链接错误只在某台机器出现，先比对 include 路径和库路径，往往比猜 C++ 语义更快。

\section{启动路径排查：从 \code{execve} 到 \code{main}}

普通 C++ 程序不是从 \code{main} 的第一行凭空开始。操作系统创建进程映像，动态链接器映射共享库，运行时建立栈和初始化环境，C/C++ runtime 调用全局构造，最后才进入 \code{main}。若问题发生在启动阶段，盯着 \code{main} 里的代码可能没有用。

启动失败可以按时间线排查。第一阶段是执行文件选择：shell 是否找到了你以为的路径，脚本是否调用了另一个可执行文件，文件是否有执行权限，架构是否匹配。第二阶段是加载：ELF 解释器是否存在，共享库是否找到，版本符号是否满足，重定位是否成功。第三阶段是运行时初始化：全局对象构造是否崩溃，线程局部对象是否初始化，环境变量是否改变库行为。第四阶段才是 \code{main} 中的用户逻辑。

\code{strace} 对启动路径非常有用，因为它能看到 \code{execve}、\code{openat}、\code{mmap}、\code{access} 等系统调用。若程序报告找不到配置文件，\code{strace} 能告诉你它实际尝试打开哪些路径。若程序报告找不到共享库，动态链接器错误、\code{ldd}、\code{readelf -d} 和环境变量要一起看。若程序在 \code{main} 前崩溃，GDB 的 backtrace、全局构造和共享库初始化顺序就进入排查范围。

\subsection{为什么启动成本不能混进 kernel 成本}

性能测量必须区分启动成本和稳态 kernel 成本。启动阶段可能包含动态链接、缺页、全局构造、内存分配、配置解析、线程池创建、JIT 或首次符号解析。这些成本可能很大，但它们不是某个热循环每次调用的成本。若把整个程序运行时间当作 kernel 时间，就会把启动噪声误判为算法差异。

这并不意味着启动成本不重要。命令行工具、短生命周期服务、插件加载和在线推理冷启动都可能关心启动时间。关键是问题陈述要清楚：你测的是端到端启动，还是核心循环吞吐。前者应保留加载和初始化；后者应把它们移出 timed region。两类实验都合理，但不能混用同一结论。

\section{从构建图理解链接错误}

链接错误常被初学者当成“编译器报错”，但它其实发生在更后面的阶段。编译每个翻译单元时，编译器只需要知道声明是否可用；链接时，链接器才需要把未定义符号解析到某个目标文件或库中的定义。若声明存在但定义缺失，编译可以通过，链接会失败。

排查链接错误要画构建图。节点包括源文件、目标文件、静态库、共享库和最终可执行文件。边表示编译或链接依赖。看到 undefined reference 时，要问：哪个目标文件引用了符号，哪个库应该提供符号，该库是否出现在链接命令中，顺序是否正确，符号名是否因为 C++ name mangling 不匹配，函数签名是否一致，目标架构是否一致，符号可见性是否隐藏。

这套方法也适用于“能链接但运行错”的情况。链接成功只表示某个符号被解析，不表示解析到的是你想要的实现。弱符号、符号 interposition、动态库搜索路径、插件系统和版本脚本都可能改变最终绑定。性能报告中若声称使用了某个优化库，应保存符号和动态加载证据，而不是只看 CMake 配置。

\section{综合案例：同一源码为什么跑出两个程序}

设想你修改了 \code{fast_sum.cpp}，本地运行 benchmark 发现没有任何变化。第一反应不应该是“优化无效”，而应该做身份审计。检查编译日志发现目标文件没有重新生成，因为构建系统没有把某个生成头文件列为依赖。清理构建目录后重新编译，二进制校验和改变，反汇编中出现新循环，benchmark 才显示变化。

另一个常见版本是动态库混用。你编译了 \filepath{build/libfast.so}，可执行文件启动时却加载 \filepath{/usr/lib/libfast.so}。源码和本地库都正确，进程执行的却是系统库。此时 \code{ldd}、\code{readelf -d}、\code{LD_DEBUG=libs} 和 \filepath{/proc/<pid>/maps} 能提供证据。若只盯着源码，会一直误判。

这个案例说明，程序不是一个抽象名词，而是一组具体产物和运行条件。底层学习从第一章就要求读者建立这种意识：任何现象都要先绑定到明确对象。后面读汇编、查 ABI、测性能，都要继承这个习惯。

\section{编译选项也是语义环境}

编译选项不只是速度开关。优化级别、目标架构、异常、RTTI、断言、浮点模型、栈保护、可见性、链接方式、LTO、调试信息，都会改变最终程序。两个可执行文件即使来自同一源码，只要编译选项不同，就应视为不同实验对象。

\code{-O0} 适合观察接近源码结构的机器码，但它不是性能对象。\code{-O3} 更接近优化构建，但会内联、删除、重排和向量化。\code{-march=native} 可能生成当前机器支持而其他机器不支持的指令。\code{-DNDEBUG} 会删除 \code{assert}，可能改变错误路径。\code{-ffast-math} 会放宽浮点语义，不能和严格浮点版本直接当作同一程序比较。

因此，报告不能只写“用 clang 编译”。要写完整命令，至少写关键 flags。若使用 CMake，也要保存 verbose build 或 compile commands。构建环境越复杂，越需要把隐含选项显式化。

\section{源文件边界和翻译单元思维}

C++ 不是把整个项目一次性解释执行。每个源文件经过预处理后形成翻译单元，编译器分别编译，再由链接器组合。这个边界影响声明可见性、内联、模板实例化、静态对象、匿名命名空间、内部链接和外部链接。

很多问题必须用翻译单元思维解释。头文件里只写声明，没有定义，编译可能通过但链接失败；在一个源文件中定义了 \code{static} 函数，另一个源文件不能引用；模板定义若不在头文件中可见，实例化可能失败；两个源文件各自有同名内部链接对象，它们不是同一个对象。链接器看到的是符号和目标文件，不是完整 C++ 语义。

优化也受边界影响。没有 LTO 时，编译器通常看不到其他翻译单元的函数体，只能按声明和 ABI 生成调用。有 LTO 时，边界可能被打通，函数被内联，符号消失，性能和反汇编都变化。看到函数“没有了”时，要先问是否内联、LTO、死代码删除或符号可见性改变，而不是立刻判断源码没有执行。

\section{运行时资源也是程序的一部分}

进程运行不仅依赖二进制，还依赖资源：配置文件、输入数据、环境变量、当前目录、文件描述符、网络、权限、CPU 亲和性、动态库、插件、临时目录。对于系统软件和性能实验，这些资源会直接改变行为。一个程序在 IDE 中能运行，在终端失败，常常只是工作目录或环境不同。

记录运行时资源，是为了让问题可复现。若程序打开相对路径，报告应写当前工作目录；若读取环境变量，报告应记录关键变量；若使用动态库和插件，报告应记录实际路径；若 benchmark 绑定 CPU 或设置线程数，报告应记录这些设置。程序身份不是二进制文件单独决定的，而是二进制加运行资源共同决定的。

\section{构建可复现性的工程底线}

可复现构建不一定要求每个字节在所有机器上完全相同，但至少要让别人知道如何得到同类产物。第一层是命令可复现：源码、编译器、参数、依赖和构建脚本明确。第二层是环境可复现：系统、容器、工具链、库版本和路径记录清楚。第三层是产物可识别：生成的目标文件和可执行文件有校验和、build id 或版本信息。第四层是运行可复现：输入、配置、环境变量和工作目录可重建。

缺少这些底线，后续章节的所有分析都会变弱。你不能稳定生成同一类二进制，就无法比较汇编；不能确认加载库，就无法讨论 ABI；不能复现输入，就无法比较性能。可复现性不是发布工程才需要的事，从第一个实验开始就应训练。

实际项目可以逐步提高要求。个人实验先保存命令和环境；团队项目加入 compile database、锁定依赖和 CI 构建；发布项目再考虑容器、构建元数据、符号服务器和制品归档。要求逐步提高，但核心思路不变：程序不是“我刚才运行的那个”，而是可命名、可追溯、可复查的产物。

\section{把程序启动写成一张时间线}

学习本章后，建议把任意一个程序的启动写成时间线。时间线从 shell 解析命令开始，到 \code{execve}，到内核加载 ELF，到账户权限和文件描述符继承，到动态链接器映射共享库，到运行时初始化，到全局构造，到进入 \code{main}，再到用户代码读取输入。每一步都写出可能失败的原因和可用工具。

这张时间线能帮助读者处理很多现实问题。程序启动时报共享库缺失，定位在动态链接阶段；全局对象构造崩溃，定位在 \code{main} 之前；相对路径找不到文件，定位在工作目录和资源阶段；第一次调用慢，可能涉及 lazy binding、缺页或初始化；benchmark 把启动时间算进去，说明计时边界错了。

时间线思维也能防止把所有问题都归咎于源码逻辑。进程是源码、二进制、运行时和操作系统共同建立的执行实例。能把启动拆开，才真正理解“程序到进程”的含义。

\section{项目里的“程序”往往不止一个}

一个仓库中常常不只有一个程序。它可能包含命令行工具、测试程序、benchmark、动态库、插件、生成器、脚本和示例。每个产物都有自己的入口、链接关系、运行资源和使用场景。读者在报告中写“程序”时，最好直接写具体产物名和路径，而不是用一个模糊名词。

这种区分对性能尤其重要。测试程序为了覆盖边界，可能使用 Debug 或 sanitizer；benchmark 为了测热路径，可能使用 Release 和固定输入；命令行工具关心端到端启动；动态库关心 ABI 和符号可见性。把这些产物混在一起，会造成错误比较。比如用测试程序的时间评价库性能，或用 benchmark 的输入假设评价用户场景，都不严谨。

构建系统也要反映这种区分。不同 target 应有清楚的编译参数、依赖和输出目录。报告中应说明分析的是哪个 target，是否带 LTO，是否静态链接，是否启用 sanitizer，是否包含测试辅助代码。只有产物身份清楚，后面的汇编、ABI 和测量才有对象。

\section{本章的最后提醒：不要跳过身份确认}

以后每次遇到底层问题，都先问“我正在分析哪个产物”。这句话看似普通，却能避免大量错误。运行路径、构建目录、动态库、输入文件、配置、环境变量、符号文件，只要有一个不匹配，后面的推理都可能偏离。

身份确认不是慢动作，而是加速器。它让你少走弯路，也让报告更可信。第一章打下的这个习惯，会一直延伸到最后一章：可信性能数据同样需要可信程序身份。

\section{一个最小身份清单}

以后做任何实验，至少写下五行：源码提交，构建命令，二进制路径，运行命令，输入路径。若涉及共享库，再加实际加载路径；若涉及性能，再加机器和编译器版本。这个清单很短，却能让多数实验从“我记得”变成“可复查”。

不要等项目变复杂才开始记录。小实验越早养成习惯，大项目越不容易混乱。

这五行还会帮助你和未来的自己沟通。几个月后重新打开实验目录时，你不需要猜当时运行了什么，也不需要凭记忆恢复环境。底层学习中，很多“遗忘”其实不是知识遗忘，而是当初没有给证据留下入口。

身份记录也会降低协作成本。别人接手你的实验时，最怕看到一堆输出文件却不知道它们来自哪次构建、哪个命令和哪个输入。清单越短，越要稳定；稳定的小清单比临时写的大段回忆更可靠。以后分析崩溃、汇编、ABI 或性能时，这个清单就是所有证据的根目录。

把身份写清楚，也是在训练工程耐心。底层问题往往不是没有答案，而是证据对象混乱导致答案无法归属。先把对象钉住，后面的推理才有落点。

这一步越早做，后面越省时间，也越少产生误判。

当实验结果需要交给别人复查时，身份清单就是最小交接协议。它让源码、构建、二进制、输入和运行环境之间的关系保持可追踪，也让后续章节中的每一条证据都有明确归属。

没有归属，证据就会失去工程价值，也难以复查。

\section{本章小结：从文本到执行实例}

现在可以把整条链路重新压缩成一句话：程序员写的是源码，编译系统把它变成带元数据的二进制文件，操作系统和运行时把二进制文件变成进程，CPU 在进程上下文中执行指令，而性能测量记录的是这个执行过程在某个环境中的表现。

这句话里的每个名词都不能省略。没有源码，就没有语言语义；没有编译系统，就没有优化和目标机器选择；没有二进制元数据，加载器不知道如何建立进程；没有操作系统和运行时，普通 C++ 程序无法获得参数、栈、库和资源；没有 CPU 执行，机器码只是文件；没有测量上下文，时间数字没有解释力。

本章收束时，不需要立刻掌握 ELF 的所有字段，也不需要背下动态链接的全部重定位类型；但必须能在脑中画出这条路径，并在遇到现象时判断应该去哪一层找原因。这是后面学习汇编、内存、ABI 和 benchmark 的共同基础。

\section{贯穿实验：拆开一个程序}

\begin{labbox}
创建文件：

\begin{lstlisting}[language=bash]
mkdir -p scratch/ch01
cat > scratch/ch01/add.cpp <<'EOF'
extern "C" int add_i32(int a, int b)
{
    return a + b;
}
EOF
\end{lstlisting}

逐层生成中间产物：

\begin{lstlisting}[language=bash]
g++ -E scratch/ch01/add.cpp -o scratch/ch01/add.i
g++ -S -O0 -masm=intel scratch/ch01/add.cpp -o scratch/ch01/add_O0.s
g++ -S -O3 -masm=intel scratch/ch01/add.cpp -o scratch/ch01/add_O3.s
g++ -c -O3 scratch/ch01/add.cpp -o scratch/ch01/add.o
objdump -drwC -Mintel scratch/ch01/add.o
\end{lstlisting}

观察任务：

\begin{enumerate}
  \item \code{.i}、\code{.s}、\code{.o} 分别是什么。
  \item \code{-O0} 和 \code{-O3} 汇编差异。
  \item \code{objdump} 输出中的机器码字节。
  \item 函数符号名是否保留为 \code{add_i32}。
\end{enumerate}
\end{labbox}

\section{实验：观察地址空间}

\begin{labbox}
\begin{lstlisting}[language=bash]
cat > scratch/ch01/address.cpp <<'EOF'
#include <iostream>
#include <vector>

int global_initialized = 42;
int global_zero;
char const* message = "hello";

int main()
{
    int stack_value = 1;
    auto* heap_value = new int(2);
    std::vector<int> vec(16, 3);

    std::cout << "main               " << reinterpret_cast<void*>(&main) << '\n';
    std::cout << "global_initialized " << &global_initialized << '\n';
    std::cout << "global_zero        " << &global_zero << '\n';
    std::cout << "message variable   " << &message << '\n';
    std::cout << "message literal    " << static_cast<void const*>(message) << '\n';
    std::cout << "stack_value        " << &stack_value << '\n';
    std::cout << "heap_value         " << heap_value << '\n';
    std::cout << "vector data        " << vec.data() << '\n';

    delete heap_value;
}
EOF

g++ -O0 -g scratch/ch01/address.cpp -o scratch/ch01/address
./scratch/ch01/address
./scratch/ch01/address
\end{lstlisting}

报告任务：

\begin{enumerate}
  \item 比较两次运行地址是否一致。
  \item 推测哪些属于 text、data、BSS、rodata、stack、heap。
  \item 解释为什么地址可能每次变化。
  \item 查看 \code{/proc/self/maps} 的方式，并解释它能提供什么信息。
\end{enumerate}
\end{labbox}

\section{把源码到进程压成因果链}

前面已经把预处理、编译、汇编、链接和加载分开讲过。真正做系统分析时，还要再往下压一层：每一层为什么必须存在，它保留什么信息，丢弃什么信息，下一层又根据这些信息作出什么决定。否则读者会知道很多工具名，却不能解释一个真实现象：源码里有一行函数调用，为什么目标文件里是一个重定位占位；链接后为什么地址变了；运行时为什么第一次访问某页才触发 page fault；调试器为什么有时看不到某个变量。

可以把 \code{add\_i32} 这个最小函数的生命周期写成一张因果账本：

\begin{center}
\begin{tabular}{p{0.18\linewidth}p{0.32\linewidth}p{0.34\linewidth}}
\toprule
层次 & 保留的信息 & 失去或改写的信息 \\
\midrule
翻译单元 & 宏展开后的完整文本、声明和定义 & 头文件边界、原始宏调用位置可能变弱 \\
AST & 名字解析、类型、表达式树、作用域 & 文本排版、注释、许多 token 细节 \\
IR & 控制流、数据流、类型化操作、副作用边界 & C++ 语法形状、部分临时对象名字 \\
优化 IR & 等价变换后的数据依赖和副作用顺序 & 未使用计算、可内联函数边界、源码局部结构 \\
汇编 & 寄存器、栈槽、指令、标签 & 高层类型、模板、对象名字的大部分语义 \\
目标文件 & 机器码字节、符号表、重定位记录、section & 最终虚拟地址、跨文件符号归属尚未确定 \\
链接后 ELF & program header、动态段、入口、最终符号绑定 & 单个目标文件的局部布局自由度 \\
进程 & 虚拟内存映射、栈、堆、共享库、线程状态 & 文件只是模板，运行状态开始独立变化 \\
\bottomrule
\end{tabular}
\end{center}

这张表的重点不是背术语，而是知道证据应该在哪一层找。若怀疑函数被内联，不能只看源码；要看优化后 IR 或反汇编。若看到目标文件里 \code{call 0x0}，不能断言程序会跳到地址 0；要看重定位。若程序第一次请求很慢，不能只看 hot loop 汇编；要看动态链接、首次触页、页缓存、全局构造和 CPU 频率。每个问题都有合适的证据层。

\subsection{AST 到 IR：优化器到底看见什么}

AST 仍然接近 C++ 语法。它知道一个表达式是二元加法，知道函数参数类型，知道名字属于哪个作用域。优化器通常不直接在完整 C++ AST 上做机器级优化，而是在更规整的 IR 上工作。IR 会把控制流、数据依赖、内存访问和调用副作用表达成更适合分析的形式。

例如：

\begin{lstlisting}[language=C++,caption={源码层的局部变量形状}]
int f(int x)
{
    int y = x + 1;
    int z = y * 2;
    return z - 2;
}
\end{lstlisting}

源码里有 \code{y} 和 \code{z}，但优化器看到的是一条可化简的数据链：\code{((x + 1) * 2) - 2} 等价于 \code{x * 2}，前提是这种变换没有违反语言语义。若有符号溢出参与，C++ 未定义行为会给优化器更大自由度；若使用无符号整数，回绕语义又会限制某些代数化简。这里第一册的语言语义和机器执行第一次连上：优化器不是按数学实数优化，而是按 C++ 抽象机器允许的行为优化。

\begin{keyidea}
编译器优化不是“猜 CPU 喜欢什么”这么简单。它先利用语言语义证明某些变换保持可观察行为，再把结果交给目标后端。C++ 语义越强，优化空间越大；程序触碰未定义行为时，优化器可以利用这个前提生成看起来反直觉的机器码。
\end{keyidea}

\subsection{目标文件：为什么重定位是核心证据}

单独编译时，编译器只能处理当前翻译单元。它知道调用了外部函数，却不知道外部函数最后放在哪里。因此目标文件必须保存两类信息：一类是已经生成的机器码和数据，另一类是“以后需要链接器修正的位置”。后一类就是重定位。

\begin{lstlisting}[numbers=none,caption={目标文件阶段的三类证据}]
nm -C caller.o
readelf -s caller.o
objdump -drwC -Mintel caller.o
\end{lstlisting}

若 \code{caller.o} 调用 \code{add\_i32}，反汇编里可能看到一个近似 \code{call 0x0} 的占位，同时 \code{objdump -r} 或 \code{readelf -r} 会显示该位置关联到符号 \code{add\_i32}。这说明机器码字节和最终地址还没有闭合。链接器读取所有目标文件和库后，才把这个位置修正为相对位移、PLT 入口或其他目标相关形式。

这个机制解释了两个常见误判。第一，目标文件反汇编不是最终运行地址，不能直接拿它解释进程中的 \register{rip}。第二，符号是否存在、调用是否经过 PLT、函数是否被静态链接或动态链接，都会改变第一次调用成本和 profile 样子。性能报告只贴源码和最终耗时，不贴符号和重定位证据，就缺少关键一层。

\subsection{execve 到 page fault：文件怎样变成运行状态}

用户在 shell 中输入命令后，关键系统调用通常是 \code{execve}。内核读取 ELF header，并不把整个文件一次性复制进内存。更常见的路径是建立虚拟内存映射：哪些文件偏移映射到进程的哪些虚拟地址，权限是只读、可执行还是可写。真正的物理页可以等到第一次访问时再通过 page fault 建立。

\begin{lstlisting}[numbers=none,caption={加载和首次触页的压缩因果链}]
execve(path)
  -> kernel validates ELF header
  -> map PT_LOAD segments into virtual address ranges
  -> map dynamic linker if needed
  -> prepare initial user stack
  -> transfer control to entry point
  -> first instruction/data access may fault
  -> kernel resolves page fault and installs page table entry
\end{lstlisting}

这条链路直接影响测量。第一次执行某段代码可能包括 I-cache 冷启动和 text page fault；第一次访问大数组可能包括匿名页分配和清零；第一次调用动态库函数可能包括 lazy binding；第一次触发 C++ 异常可能加载 unwind 相关路径。若 benchmark 把这些成本混进 hot loop，就会把 loader、VM 和 runtime 成本误判成算法成本。

\subsection{系统边界：syscall、异常和中断不是普通函数调用}

用户态函数调用遵循 ABI：参数寄存器、返回寄存器、caller-saved/callee-saved 和栈对齐。系统调用跨越的是权限边界。用户态通过专门指令进入内核，CPU 切换到更高权限级别，内核根据 syscall number 分派处理，完成后再返回用户态。这个路径涉及寄存器约定、内核栈、权限检查、可能的调度和安全隔离，不能用普通函数调用成本类比。

page fault 也不是普通错误返回。CPU 执行 load/store 或取指时发现页表项缺失、权限不允许或地址非法，会进入异常处理。内核判断这是合法按需分配、文件映射缺页、写时复制，还是非法访问。如果合法，内核修复页表并回到触发异常的指令重新执行；如果非法，向进程发送信号。于是同一条 \instruction{mov} 指令，可能在热路径中只是 L1 load，也可能第一次触页时跨入内核。

\begin{center}
\begin{tabular}{p{0.22\linewidth}p{0.34\linewidth}p{0.30\linewidth}}
\toprule
边界 & 触发原因 & 性能和调试含义 \\
\midrule
普通函数调用 & \instruction{call} 到用户态函数 & ABI、栈帧、内联和 PLT 影响成本 \\
系统调用 & 用户态请求内核服务 & 权限切换、内核路径、可能阻塞 \\
page fault & 虚拟地址翻译失败或权限不符 & 可能是正常按需分页，也可能是 bug \\
硬件中断 & 设备或定时器事件 & 可打断用户线程，影响延迟分布 \\
信号 & 内核向进程投递事件 & 控制流异步改变，调试要看上下文 \\
\bottomrule
\end{tabular}
\end{center}

成熟的底层报告必须区分这些边界。若一次请求 p99 偶发升高，可能不是 C++ 函数突然慢，而是 page fault、调度、I/O、中断或系统调用长尾进入了计时范围。

\section{展开：优化器如何在 IR 上证明“可以改写”}

只说“编译器会优化”还不够。经典系统教材会继续追问：优化器到底看见什么，凭什么知道某个 load 可以提前，凭什么知道某个分支可以删掉，凭什么把一个函数调用消成常量。答案不在源码文本里，而在 IR 中的控制流、数据流、类型、别名和副作用模型。

\subsection{从表达式树到 SSA：变量名不再是核心}

源码里的变量名服务程序员阅读；优化器更关心“哪个定义产生了哪个值”。许多现代 IR 使用 SSA 形态：每个虚拟值只定义一次，控制流汇合处用 \code{phi} 选择来自不同前驱块的值。下面这个例子很小，但足以说明优化器的视角。

\begin{lstlisting}[language=C++,caption={源码中的变量会反复赋值}]
int clamp_add(int x, int y)
{
    int z = x + y;
    if (z < 0) {
        z = 0;
    }
    return z + 1;
}
\end{lstlisting}

对应的 SSA 直觉可以写成：

\begin{lstlisting}[numbers=none,caption={SSA 形态关注定义和控制流}]
entry:
  z0 = add x, y
  c0 = icmp_slt z0, 0
  branch c0, neg, pos

neg:
  z1 = 0
  jump merge

pos:
  z2 = z0
  jump merge

merge:
  z3 = phi [z1 from neg], [z2 from pos]
  r0 = add z3, 1
  return r0
\end{lstlisting}

\code{phi} 不是机器指令，它是 IR 层表达“根据控制流来源选择值”的记账方式。后端会把它消成寄存器移动、条件移动或分支布局。理解 SSA 后，再看优化就不神秘：常量传播是在值图上传播常量，死代码删除是删掉没有使用且没有副作用的定义，公共子表达式消除是复用等价定义，循环优化是在控制流和依赖图上改写迭代结构。

\subsection{优化 pass 的因果链}

一个 pass 不能随便改程序。它必须有前置事实、改写规则和保持的语义。下面的表把常见优化压成可检查形式：

\begin{center}
\begin{tabular}{p{0.22\linewidth}p{0.32\linewidth}p{0.32\linewidth}}
\toprule
优化 & 需要证明什么 & 容易被谁阻止 \\
\midrule
常量传播 & 某个 SSA 值在所有到达路径上相同 & 未知函数调用、输入、volatile \\
死代码删除 & 定义结果没人用，且无可观察副作用 & store、I/O、atomic、异常边 \\
内联 & 调用目标可见，收益大于成本 & 跨翻译单元、可见性、递归、代码体积 \\
循环不变外提 & 表达式每次迭代相同且提前不改变异常/副作用 & 可能别名的 store、越界、除零、调用 \\
向量化 & 迭代间独立，内存访问可描述，尾部可处理 & 数据依赖、别名、非规则控制流 \\
\bottomrule
\end{tabular}
\end{center}

这张表解释了为什么 C++ 语义会直接影响性能。若指针可能别名，优化器不能轻易把 load 外提；若表达式可能触发未定义行为，优化器会基于“合法程序不会触发它”的前提改写；若函数可能抛异常或访问全局状态，很多局部改写会被阻止。性能工程不能把优化器当魔法，也不能把它当敌人；要给它足够清楚的语义合同。

\subsection{一个反例：源码短不代表机器路径短}

下面两个函数在源码里都很小，但优化空间完全不同。

\begin{lstlisting}[language=C++,caption={两个看似相近的求和接口}]
int sum_plain(int const* p, int n)
{
    int s = 0;
    for (int i = 0; i < n; ++i) {
        s += p[i];
    }
    return s;
}

int sum_with_call(int const* p, int n);

int caller(int const* p, int n)
{
    return sum_with_call(p, n) + 1;
}
\end{lstlisting}

若 \code{sum_plain} 在当前翻译单元可见，优化器可能展开循环、做多累加器、甚至在目标允许时向量化。若 \code{sum_with_call} 只有声明，普通编译阶段不知道函数体，只能生成调用。启用 LTO 后，IR 可能跨翻译单元进入链接阶段，优化器才有机会内联和继续传播。报告中若只写“函数调用很慢”或“编译器没优化”，是不完整的；应写明是否同翻译单元、是否启用 LTO、符号可见性、是否经过 PLT、反汇编里是否真的存在调用。

\section{展开：ELF、动态链接和异常展开的运行时账本}

ELF 不是文件格式知识点，而是程序能被加载、链接、定位、调试、展开栈和安全隔离的合同。只看 \code{.text} 机器码，会漏掉很多性能和故障根因。

\subsection{section 和 segment 解决的是不同问题}

节 section 主要服务链接器和分析工具；段 segment 主要服务加载器。目标文件里，\code{.text}、\code{.rodata}、\code{.data}、\code{.bss}、\code{.symtab}、\code{.rela.text} 等 section 帮助链接器放置符号和修正重定位。可执行文件运行时，内核主要看 program header 中的 \code{PT\_LOAD} 段，把文件范围映射成虚拟内存范围，并设置读、写、执行权限。

\begin{center}
\begin{tabular}{p{0.20\linewidth}p{0.34\linewidth}p{0.32\linewidth}}
\toprule
对象 & 主要使用者 & 典型问题 \\
\midrule
section & 编译器、汇编器、链接器、调试工具 & 符号在哪，重定位在哪，调试信息在哪 \\
segment & 内核 loader、动态链接器 & 哪些页映射到进程，权限是什么 \\
symbol & 链接器、动态链接器、调试器 & 名字如何绑定到地址或偏移 \\
relocation & 链接器、动态链接器 & 哪些机器码或数据位置需要修正 \\
\bottomrule
\end{tabular}
\end{center}

一个 stripped 可执行文件可能没有完整 section header，仍然能运行，因为加载器靠 segment；但人类分析会困难，因为符号和调试信息少了。反过来，目标文件有丰富 section，却不能直接运行，因为它还没有最终地址和完整加载合同。

\subsection{PLT/GOT 和 lazy binding 的第一次调用成本}

动态链接时，调用外部函数常通过 PLT 和 GOT 完成。PLT 是一段跳板代码，GOT 是保存真实地址或待解析入口的表。若启用 lazy binding，第一次调用某个外部函数时，PLT 可能先进入动态链接器解析符号，把结果写入 GOT；后续调用再直接跳到真实地址。

\begin{lstlisting}[numbers=none,caption={动态函数调用的压缩路径}]
first call:
  caller -> PLT entry
         -> dynamic linker resolver
         -> lookup symbol in loaded objects
         -> patch GOT slot
         -> real function

later call:
  caller -> PLT entry
         -> GOT slot
         -> real function
\end{lstlisting}

这解释了为什么 benchmark 要预热，为什么 \code{-fno-plt}、静态链接、\code{LD\_BIND\_NOW}、符号可见性和链接选项会改变 profile。并不是每个项目都要避免 PLT；更重要的是报告要知道调用路径是什么。若热循环里每次调用一个外部小函数，间接跳转、不可内联和符号绑定都会进入性能账本。

\subsection{TLS、初始化数组和 main 之前的工作}

\code{main} 之前，运行时还可能执行很多工作：动态链接器映射共享库，处理重定位，初始化 TLS，执行 \code{.preinit\_array}、\code{.init\_array} 和 C++ 全局对象构造，设置标准库状态和线程局部对象。线程局部变量也不是普通全局变量；不同 TLS 模型会影响访问方式，动态库中的 TLS 可能比本地 exec 模型更复杂。

\begin{lstlisting}[numbers=none,caption={main 之前可能发生的工作}]
kernel transfers control to entry
runtime startup sets stack and libc state
dynamic linker maps shared libraries
relocations are applied
TLS blocks are prepared
init arrays and global constructors run
main(argc, argv, envp) is called
\end{lstlisting}

这条链路对 AI runtime 尤其重要。模型服务启动时可能 mmap 权重、初始化线程池、创建 backend context、构造 tokenizer、注册算子和预热 kernel。若把这些成本混进首个请求 TTFT，就会把初始化成本误判成推理成本。服务报告必须拆开 process startup、model load、warmup、prefill、decode 和 streaming。

\subsection{异常展开和调试信息也是运行时合同}

C++ 异常不是简单返回错误码。抛出异常时，运行时要根据 unwind 信息沿栈查找 handler，执行离开作用域对象的析构。ELF 中的 unwind 表、语言运行时、ABI 和编译器生成的清理路径共同决定异常能否正确传播。即使项目禁止热路径异常，也要理解这层：析构、RAII、栈展开、core dump 和 profiler 栈回溯都依赖类似元数据。

\begin{center}
\begin{tabular}{p{0.24\linewidth}p{0.34\linewidth}p{0.28\linewidth}}
\toprule
机制 & 作用 & 性能/调试影响 \\
\midrule
unwind table & 描述如何从当前 PC 恢复调用者上下文 & 栈回溯、异常传播、profile 采样 \\
landing pad & 异常路径上的清理和 handler 入口 & RAII 析构、异常边界 \\
debug info & 地址到源码、变量和类型的映射 & 调试、符号化、崩溃分析 \\
build-id & 绑定二进制与符号文件 & core dump 与正确符号匹配 \\
\bottomrule
\end{tabular}
\end{center}

生产系统中，符号、build-id、unwind 信息和编译选项是事故分析材料的一部分。没有这些材料，性能火焰图可能只剩地址，崩溃栈可能断掉，异常路径可能被误读。第一册讲 ELF 的目的不是背字段，而是让读者知道这些“元数据”会决定真实工程能否诊断。

\section{状态机：从 \code{execve} 到第一条业务指令}

经典教材厚，不是因为把术语重复很多遍，而是把“程序开始运行”拆成可验证状态。下面这台状态机可以作为所有启动、冷启动、动态链接和首个请求分析的底图。

\begin{center}
\begin{tabular}{p{0.22\linewidth}p{0.30\linewidth}p{0.34\linewidth}}
\toprule
状态 & 进入条件 & 可观察证据 \\
\midrule
FileChosen & shell 或父进程给出路径和 argv/envp & 命令行、工作目录、环境变量 \\
ElfValidated & 内核读取 ELF header 并检查格式 & \cmd{readelf -h}、错误码 \\
SegmentsMapped & 内核建立 \code{PT\_LOAD} 虚拟映射 & \filepath{/proc/self/maps}、权限位 \\
InterpreterMapped & 动态链接器被映射并取得控制权 & \cmd{readelf -l} 的 interpreter \\
RelocationPending & 共享库和重定位表已知但未完全解析 & \cmd{readelf -d}、\cmd{readelf -r} \\
RuntimeInit & TLS、init array、全局构造和 libc 状态初始化 & 构造日志、\cmd{LD\_DEBUG}、trace \\
MainEntered & 用户 \code{main} 获得控制权 & 断点、日志、采样栈 \\
HotLoopReady & 输入和工作区准备完毕，进入被测区域 & benchmark harness 的 timed region \\
\bottomrule
\end{tabular}
\end{center}

这张表的用处是排除错误归因。若首个请求慢发生在 \code{RuntimeInit} 或 \code{RelocationPending}，就不能把它算成模型 decode 慢；若慢发生在 \code{HotLoopReady} 之后，再去看汇编、cache 和 PMU 才有意义。状态机也能解释为什么同一个二进制在两台机器上冷启动差异很大：动态库路径、页缓存、ASLR、符号绑定、CPU 频率和全局构造都可能不同。

\subsection{动态绑定的状态机}

PLT/GOT 的第一次调用也可以写成状态机。这样比一句“第一次会慢”更可审查。

\begin{lstlisting}[numbers=none,caption={lazy binding 状态机}]
Unresolved:
  GOT slot points to resolver trampoline
  first call enters PLT

Resolving:
  dynamic linker searches loaded objects
  symbol version and visibility are checked
  relocation target is computed

Patched:
  GOT slot is updated to real function address
  later calls skip resolver

Bound:
  hot path contains normal indirect jump/call through GOT
\end{lstlisting}

对应反例也很重要。若设置 \code{LD\_BIND\_NOW=1}，许多符号会在启动阶段解析，首个业务调用不再承担 lazy binding；若静态链接，路径又不同；若启用 \code{-fno-plt}，调用形态也会变化。报告不能只写“第一次慢”，而要写动态绑定策略和观测证据。

\subsection{page fault 的状态机}

page fault 不是一种单一错误，而是虚拟内存状态机中的一次转移。

\begin{center}
\begin{tabular}{p{0.24\linewidth}p{0.34\linewidth}p{0.28\linewidth}}
\toprule
状态 & 触发 & 后果 \\
\midrule
MappedNoPte & VMA 存在，但页表项尚未建立 & 首次访问进入内核 \\
MinorFault & 物理页已在内存或可零页建立 & 更新页表，重试指令 \\
MajorFault & 需要从文件或设备读入页 & I/O 进入延迟账本 \\
CowFault & 写共享只读页 & 复制物理页，再重试 store \\
ProtectionFault & 权限不符或非法地址 & 信号、崩溃或异常处理路径 \\
\bottomrule
\end{tabular}
\end{center}

性能后果可以用一个简单实验验证：同一块 1GiB 匿名内存，第一次顺序写包含缺页和清零，第二次顺序写更多反映内存带宽；若每次 benchmark 都重新分配并首次触碰，就测到了 VM 成本，不是纯粹 store 吞吐。工程事故常发生在这里：服务上线后首个大请求 p99 高，被误判为模型 kernel 慢，实际是权重 mmap 后首次触页、透明大页整理或 COW 路径进入了计时范围。

\subsection{实验：把启动路径变成证据}

\begin{lstlisting}[numbers=none,caption={启动路径实验命令}]
readelf -h ./app
readelf -l ./app
readelf -d ./app
readelf -r ./app | head
LD_DEBUG=libs,bindings ./app 2>ld-debug.txt
cat /proc/$(pgrep app)/maps
\end{lstlisting}

实验报告要回答五个问题：入口 interpreter 是谁；哪些共享库被加载；是否存在 lazy binding；\code{main} 前是否有全局构造日志；第一次业务调用是否和第二次调用处于同一状态。能回答这些问题，才算把“从文件到进程”的状态机真正跑通。

\subsection{证据包：真实 ELF、符号和反汇编}

本章的验收不能只停在“知道有 ELF”。下面是一台 AArch64 Linux 环境上对一个最小程序的真实证据片段。程序中保留一个 \code{noinline} 的 \code{add\_noinline}，目的是让符号和反汇编中有可追踪目标。

\begin{lstlisting}[numbers=none,caption={最小证据命令}]
g++ -O2 -g tiny.cpp -o tiny
readelf -h tiny
readelf -l tiny
readelf -d tiny
nm -C tiny
objdump -d --demangle tiny
perf stat -e cycles,instructions ./tiny
\end{lstlisting}

真实输出中最关键的字段如下：

\begin{lstlisting}[numbers=none,caption={ELF 和动态链接证据片段}]
ELF Header:
  Class:   ELF64
  Data:    2's complement, little endian
  Type:    EXEC (Executable file)
  Machine: AArch64
  Entry point address: 0x400640

INTERP:
  [Requesting program interpreter: /lib/ld-linux-aarch64.so.1]

Dynamic section:
  NEEDED Shared library: [libstdc++.so.6]
  NEEDED Shared library: [libm.so.6]
  NEEDED Shared library: [libgcc_s.so.1]
  NEEDED Shared library: [libc.so.6]
\end{lstlisting}

这些字段把状态机落到文件事实：\code{Entry point} 不是 \code{main}；\code{INTERP} 说明先由动态加载器接手；\code{NEEDED} 说明启动路径包含共享库解析。若冷启动 benchmark 没记录这些字段，就无法区分“业务 kernel 慢”和“加载/绑定/初始化慢”。

\begin{lstlisting}[numbers=none,caption={符号和反汇编证据片段}]
nm -C tiny:
  0000000000400780 T add_noinline
  00000000004005c0 T main

objdump:
  4005e8: 94000066  bl  400780 <add_noinline>

0000000000400780 <add_noinline>:
  400780: 0b010000  add w0, w0, w1
  400784: d65f03c0  ret
\end{lstlisting}

这里可以把 C++、ABI 和机器指令连起来：AArch64 的整数参数在 \code{w0/w1}，返回值在 \code{w0}，\code{bl} 写入返回地址并跳转，\code{ret} 返回调用点。若优化器把函数 inline 掉，\code{nm} 和 \code{objdump} 会出现不同证据；这不是工具坏了，而是优化状态变了。

\subsection{PMU 不可用也是证据}

同一环境中，\cmd{perf stat} 可能返回：

\begin{lstlisting}[numbers=none,caption={PMU 权限或平台受限时的真实形态}]
Error:
No supported events found.
Access to performance monitoring and observability operations is limited.
\end{lstlisting}

这时报告不能编造 cycles、instructions 或 cache miss。正确降级是保存 \code{readelf}/\code{objdump}、输入 checksum、阶段时间、循环次数、机器信息和源码版本；结论也要写窄：可以说“反汇编显示热循环确实调用了目标函数”，不能说“IPC 证明瓶颈在前端”。经典系统教材强调证据边界，原因就在这里。

\section{编译器和运行时深水区：从规则到事故}

第一册如果只讲“源码变成汇编”，仍然不够。真实系统事故经常发生在更深的连接处：C++ 抽象机给优化器提供了什么前提；IR 如何保存别名、异常和副作用；ELF 元数据如何影响加载、符号解析和回溯；运行时库、分配器、线程库和内核边界如何改变一次请求的成本。本节把这些深水区压成一张地图，后续章节和实验可以逐项展开。

\subsection{优化器深水区：不是所有源码事实都会留到机器码}

优化器做任何改写，都依赖某种可证明事实。问题在于，很多事实不是程序员的直觉，而是 C++ 标准、编译器 IR 和目标平台共同给出的合同。下面的表把常见深水区和证据层对齐：

\begin{center}
\small
\begin{tabular}{p{0.22\linewidth}p{0.30\linewidth}p{0.32\linewidth}}
\toprule
深水区 & 优化器可能利用的前提 & 应保留的证据 \\
\midrule
signed overflow & 合法 C++ 程序不会让有符号整数溢出 & UBSan、优化 IR、\code{-O0/-O3} 汇编对照 \\
strict aliasing & 不相容类型的对象通常不会互相别名 & 源码类型、IR alias metadata、错误强转反例 \\
object lifetime & 只有 alive object 可以按其类型访问 & placement new、\code{std::launder}、ASan/UBSan 报告 \\
escape analysis & 未逃逸对象可放入寄存器或消除分配 & IR alloc 消除前后、反汇编、分配计数 \\
exception edge & 可能抛异常的表达式限制移动和删除 & landing pad、unwind table、优化 remark \\
atomic/volatile & 可观察副作用限制重排和删除 & IR memory order、反汇编、线程测试 \\
\bottomrule
\end{tabular}
\end{center}

例如一个越界访问在 Debug 下似乎只读到相邻内存，Release 下却把整个分支删掉；这不应该先解释成“编译器 bug”。正确路径是：确认源码是否有未定义行为，打开 sanitizer，查看优化前后 IR，再看汇编是否基于“合法程序不会走到这里”的前提改写。只有这条链闭合，才有资格讨论编译器是否违反语义。

\begin{lstlisting}[numbers=none,caption={优化器证据包的最小命令}]
clang++ -O0 -g -fsanitize=address,undefined bug.cpp -o bug_asan
clang++ -O3 -S -emit-llvm bug.cpp -o bug_O3.ll
clang++ -O3 -S -masm=intel bug.cpp -o bug_O3.s
clang++ -O3 -Rpass=loop-vectorize bug.cpp -c -o bug.o
objdump -drwC -Mintel bug.o
\end{lstlisting}

报告中要把“语言错误”“优化事实”和“机器路径”分开写。Sanitizer 抓到越界，只能证明程序违反合同；IR 说明优化器看到什么；反汇编说明最终执行什么。把三者混成一句“Release 崩了”没有工程价值。

\subsection{运行时深水区：main 之外还有很多状态机}

运行时不是一个模糊词。它包含动态链接器、C/C++ 运行时、异常运行时、分配器、线程库、系统调用包装、vDSO、信号处理和内核对象。它们共同决定程序进入 \code{main} 前、热路径外和故障路径上的行为。

\begin{center}
\small
\begin{tabular}{p{0.22\linewidth}p{0.34\linewidth}p{0.28\linewidth}}
\toprule
运行时边界 & 它负责什么 & 典型误判 \\
\midrule
dynamic linker & 映射共享库、处理重定位、符号绑定 & 把首次 lazy binding 算成 kernel 时间 \\
TLS runtime & 为每个线程准备线程局部存储 & 把 TLS 访问当普通全局 load \\
exception runtime & 查找 handler、执行析构、栈展开 & 崩溃栈缺失时只怪 GDB \\
allocator & 管理 heap、arena、mmap、缓存和回收 & 把分配器锁竞争误判为算法慢 \\
pthread runtime & 创建线程、阻塞、唤醒、join、TLS 清理 & 把线程启动成本混进 steady state \\
syscall/vDSO & 跨内核边界或走用户态快捷页 & 把 \code{clock_gettime} 成本当普通函数调用 \\
signal/page fault & 异步控制流、缺页和权限异常 & 把偶发 p99 都归咎于 C++ 代码 \\
\bottomrule
\end{tabular}
\end{center}

这些边界会直接影响后两册。第二册讨论线程池和同步时，\code{pthread}、futex、调度和 TLS 析构都是运行时账本；第三册讨论模型服务时，权重 \code{mmap}、首次触页、allocator arena、线程池启动、backend context 和 tokenizer 初始化都会进入 TTFT。若报告只写 \code{main} 里的函数耗时，就无法解释生产中的冷启动和尾延迟。

\subsection{深水区报告字段}

遇到 Release-only 崩溃、冷启动慢、第一次请求慢、栈回溯断裂、函数在符号表中消失、benchmark 数字异常时，报告至少应包含下面字段：

\begin{lstlisting}[numbers=none,caption={CompilerRuntimeDeepReport}]
CompilerRuntimeDeepReport:
  source_commit
  compiler_and_flags
  sanitizer_result
  optimization_ir_or_remark
  object_symbols_and_relocations
  final_binary_headers
  dynamic_linking_mode
  startup_state:
    loader
    relocations
    init_array
    tls
  runtime_state:
    allocator
    thread_count
    syscall_or_vdso_path
    page_fault_count
  machine_evidence:
    objdump
    maps
    backtrace_or_unwind
    timing_scope
\end{lstlisting}

这份报告不是为了繁琐，而是为了避免错误层级。若符号在目标文件中存在、链接后被内联消失，这是编译/链接问题；若二进制正确但运行时加载了旧共享库，这是部署问题；若 warmup 后速度正常但首个请求慢，这是加载、触页或初始化问题；若 sanitizer 失败，先修定义性问题，再谈性能。

\subsection{最小验收：一个问题至少跨三层证明}

本册后续每次解释一个底层问题，至少要跨三层。解释优化，必须有 C++ 规则、IR 或优化 remark、反汇编；解释链接，必须有符号、重定位、最终 ELF；解释运行时，必须有 \code{/proc/<pid>/maps}、动态库或 syscall/page fault 证据；解释崩溃，必须有寄存器、地址、对象生命周期和符号化栈。

这样做会让章节变厚，但厚度来自必要证据。大师级系统教材不满足于“知道一个概念”，而是训练读者在源码、二进制、运行时和机器状态之间来回闭合。

\section{六条必须打穿的深水链路}

前面已经建立了源码到进程的大地图。现在把第一册仍然容易薄的部分继续压深：优化器、IR、ELF 运行时元数据、DWARF 调试信息、C++ 异常 ABI 和系统调用边界。它们不是孤立高级主题，而是同一条问题链的不同断面：一段 C++ 为什么会变成这段机器行为，工具为什么这样显示，运行时为什么在这里花时间或崩溃。

\subsection{一、从 AST 到 IR、pass 和汇编的闭环}

只看最终汇编，读者容易把优化器当成黑箱；只看源码，又无法解释机器行为。合格的第一册应该让同一个函数依次经过 AST、IR、优化 pass 和汇编。下面固定一个小函数：

\begin{lstlisting}[language=C++,caption={用于观察优化链路的小函数}]
int sum_positive(int const* p, int n)
{
    int sum = 0;
    for (int i = 0; i < n; ++i) {
        if (p[i] > 0) {
            sum += p[i];
        }
    }
    return sum;
}
\end{lstlisting}

源码层能看到变量、循环和 \code{if}；AST 能看到语法树和类型；未优化 IR 往往还能看到栈槽、\code{load} 和 \code{store}；优化后 IR 更接近 SSA 数据流；后端再把 SSA 值分配到寄存器、生成分支或条件移动，最后才形成汇编。观察命令可以固定为：

\begin{lstlisting}[numbers=none,caption={同一函数的编译器证据链}]
clang++ -O0 -Xclang -ast-dump -fsyntax-only sum.cpp > ast.txt
clang++ -O0 -S -emit-llvm sum.cpp -o sum_O0.ll
clang++ -O3 -S -emit-llvm sum.cpp -o sum_O3.ll
clang++ -O3 -S -masm=intel sum.cpp -o sum_O3.s
clang++ -O3 -Rpass=loop-vectorize -Rpass-missed=loop-vectorize \
  -c sum.cpp -o sum.o 2>remarks.txt
objdump -drwC -Mintel sum.o > sum.objdump.txt
\end{lstlisting}

这组证据回答不同问题。AST 说明编译器前端理解到什么语法和类型；\code{sum\_O0.ll} 说明源代码被降低成怎样的内存和控制流；\code{sum\_O3.ll} 说明优化器证明了哪些值、分支和内存访问关系；\code{remarks.txt} 说明循环是否向量化以及为什么失败；反汇编说明最终机器路径。只有这五层对齐，才能说“我理解这段 C++ 如何执行”。

\subsection{二、LLVM IR/GCC GIMPLE：SSA、load/store 和 pass 的真实对象}

LLVM IR 和 GCC GIMPLE 都不是汇编的漂亮写法，而是优化器工作的对象。IR 的重点不是保留源码外观，而是显式表示值、控制流、内存访问和副作用。未优化 LLVM IR 中常见 \code{alloca/load/store}，因为它先把局部变量保守地放进栈槽；\term{mem2reg}{memory to register promotion} 会把可证明没有地址逃逸的栈槽提升成 SSA 值。

\begin{lstlisting}[numbers=none,caption={mem2reg 前后的直觉}]
before mem2reg:
  %sum.addr = alloca i32
  store i32 0, ptr %sum.addr
  %old = load i32, ptr %sum.addr
  %new = add i32 %old, %x
  store i32 %new, ptr %sum.addr

after mem2reg:
  %sum.0 = phi [0, %entry], [%sum.1, %loop.latch]
  %sum.1 = add i32 %sum.0, %x
\end{lstlisting}

\code{phi} 不是机器指令，而是控制流汇合处的值选择。后端会把它消成寄存器移动、分支布局、条件移动或别的机器实现。理解这一点后，调试优化代码里“变量消失”就不奇怪：源码变量可能已经不再对应一个稳定栈地址，而是分裂成多个 SSA 值。

几个核心 pass 要能从前提、改写和证据三方面解释：

\begin{center}
\small
\begin{tabular}{p{0.18\linewidth}p{0.31\linewidth}p{0.35\linewidth}}
\toprule
pass & 它证明什么 & 证据应该看哪里 \\
\midrule
mem2reg & 栈槽没有复杂地址逃逸，可提升成 SSA 值 & \code{alloca/load/store} 是否减少，是否出现 \code{phi} \\
inline & 调用目标可见且收益足够 & IR 中调用是否消失，符号表中函数是否仍独立存在 \\
DCE & 计算结果无使用且无可观察副作用 & 无用循环或临时值是否被删除 \\
LICM & 表达式循环内不变，提前不改变异常和副作用 & load 是否被外提，别名和调用是否阻止外提 \\
vectorize & 迭代间独立，内存访问可描述，尾部可处理 & remark、向量 IR、\instruction{vmov}/\instruction{vadd} 等指令 \\
\bottomrule
\end{tabular}
\end{center}

GCC 的 GIMPLE 也服务同类目标，只是术语和 dump 形态不同。可以用：

\begin{lstlisting}[numbers=none,caption={GCC 观察优化阶段的入口}]
g++ -O3 -fdump-tree-gimple -fdump-tree-optimized sum.cpp -c
g++ -O3 -fopt-info-vec-optimized -fopt-info-vec-missed sum.cpp -c
\end{lstlisting}

不要把 LLVM IR 或 GIMPLE 当作跨版本稳定教材语法背诵。它们的具体输出会随编译器版本变化。教学重点是：优化器不是按源码行工作，而是按控制流图、SSA 值、内存依赖、别名事实、异常边和副作用事实工作。

\subsection{三、ELF 深水区：\code{.init\_array}、TLS、可见性和动态链接器}

ELF 深水区不止 \code{.text}、符号表和 PLT/GOT。很多“main 之前”和“第一次调用”的问题来自 \code{.init\_array}、TLS、符号可见性、动态段和动态链接器策略。

\code{.init\_array} 保存运行时进入 \code{main} 前要调用的初始化函数地址，C++ 全局对象构造、某些库初始化和编译器生成的启动钩子都可能进入这里。一个 benchmark 若在全局对象构造里分配大内存、打开文件、注册插件或初始化日志，整个进程时间就已经包含这些成本。

\begin{lstlisting}[numbers=none,caption={观察 init array 和动态段}]
readelf -S ./app | rg 'init_array|fini_array|dynamic|dynsym'
readelf -x .init_array ./app
readelf -d ./app
objdump -drwC -Mintel ./app | rg '@plt|init_array|__libc_start'
\end{lstlisting}

TLS 也不是普通全局变量。线程局部对象需要每个线程有自己的存储实例；不同 TLS 模型会生成不同访问序列。主可执行文件中本地 exec 模型可能很轻；动态库中的 general dynamic 模型可能要调用 TLS helper。若热路径频繁访问 \code{thread\_local} 大对象或复杂对象，不能只按普通全局 load 理解。

符号可见性决定动态链接边界，也会影响优化。默认可见的共享库符号可能被外部引用、插桩或替换，编译器和链接器就要保留更强 ABI 边界；隐藏符号给优化器和链接器更多自由。库工程里常见：

\begin{lstlisting}[numbers=none,caption={符号可见性证据}]
readelf -Ws libfoo.so | rg 'GLOBAL|LOCAL|HIDDEN|DEFAULT'
objdump -T libfoo.so | rg foo
clang++ -fvisibility=hidden -fvisibility-inlines-hidden ...
\end{lstlisting}

动态链接器的证据不能只靠猜。若怀疑首次调用慢，可以打开：

\begin{lstlisting}[numbers=none,caption={动态链接器运行时证据}]
LD_DEBUG=libs,bindings ./app 2>ld-debug.txt
LD_BIND_NOW=1 ./app
readelf -d ./app | rg 'NEEDED|RUNPATH|RPATH|BIND_NOW|FLAGS'
\end{lstlisting}

\code{LD\_BIND\_NOW=1} 把许多 lazy binding 提前到加载阶段，能帮助区分“第一次业务调用慢”到底来自动态绑定，还是来自页面首次触碰、cache 冷启动或真正算法路径。符号可见性、RPATH/RUNPATH、\code{NEEDED} 版本和加载顺序，都是生产事故材料，不是文件格式细枝末节。

\subsection{四、DWARF、CFI 和 optimized out：工具为什么看不到变量}

DWARF 调试信息把机器地址、源码行、变量、类型和栈展开规则联系起来。它不是 CPU 执行程序所需的东西，却决定 GDB、crash dump、profiler 和异常展开能否把机器状态解释回人能理解的调用栈。

调试信息至少包含几类重要事实：

\begin{center}
\small
\begin{tabular}{p{0.20\linewidth}p{0.36\linewidth}p{0.28\linewidth}}
\toprule
DWARF/展开事实 & 它说明什么 & 典型工具 \\
\midrule
line table & 指令地址和源码行的大致关系 & \cmd{llvm-dwarfdump} \\
location list & 变量在某段 PC 范围内位于哪里 & GDB、\cmd{llvm-dwarfdump} \\
CFI/\code{.eh\_frame} & 如何恢复调用者 CFA、返回地址和保存寄存器 & \cmd{readelf} \\
type DIE & 类型、字段、作用域和模板实例信息 & GDB、\cmd{llvm-dwarfdump} \\
build-id & 二进制和符号文件的身份绑定 & core dump、符号服务器 \\
\bottomrule
\end{tabular}
\end{center}

\code{optimized out} 通常不是调试器偷懒，而是变量没有稳定位置。一个源码变量可能在入口处位于 \register{edi}，循环中变成某个 SSA 值，后来被常量传播，最终完全没有机器位置。调试器只能按 DWARF location list 在某个 PC 范围内恢复它；若该范围没有位置，就显示 optimized out。

\begin{lstlisting}[numbers=none,caption={调试信息和展开信息证据}]
clang++ -O2 -g f.cpp -o f
llvm-dwarfdump --debug-line f | head
llvm-dwarfdump --debug-info f | rg 'DW_TAG_variable|DW_AT_location'
readelf --debug-dump=frames f | head -80
objdump --dwarf=decodedline f | head
\end{lstlisting}

CFI 的核心概念是 \term{CFA}{canonical frame address}。它告诉展开器：在当前指令位置，调用者栈帧的基准地址如何计算，返回地址在哪里，哪些 callee-saved 寄存器保存在哪里。若编译器省略帧指针，CFI 更重要；若手写汇编没有正确 CFI，GDB 回溯、perf 采样栈和异常穿越该帧都可能失败。

\subsection{五、C++ exception ABI：throw 不是长跳转这么简单}

C++ 异常机制跨越语言语义、ABI、运行时库和 unwind metadata。一个 \code{throw} 通常会构造异常对象，调用运行时入口，例如 Itanium C++ ABI 系列实现中的 \code{\_\_cxa\_throw}；展开器沿栈查找 handler；每一帧的 personality routine 和 LSDA 描述该帧有哪些 landing pad、catch 类型和 cleanup；离开作用域的对象析构通过 cleanup 路径执行。

\begin{lstlisting}[numbers=none,caption={异常传播的压缩状态机}]
ThrowCreated:
  exception object and type info are prepared

SearchPhase:
  unwinder walks frames and asks personality routines
  whether a handler exists

CleanupPhase:
  frames are unwound
  destructors and cleanup landing pads run

HandlerEntered:
  matching catch receives control

Unhandled:
  terminate is called if no handler is found
\end{lstlisting}

这条链解释 RAII 为什么能在异常中成立：不是异常自动懂资源，而是编译器为作用域对象生成了 cleanup 边，异常运行时按 unwind 表穿过这些边。也解释为什么裸汇编和 C ABI 边界危险：如果栈帧没有可展开信息，或者让 C++ 异常穿过不理解异常 ABI 的边界，析构和回溯都可能出错。

异常证据可以这样采集：

\begin{lstlisting}[numbers=none,caption={异常 ABI 证据入口}]
clang++ -O2 -g -fno-omit-frame-pointer throw.cpp -o throw_app
nm -C throw_app | rg '__cxa_throw|personality|Unwind'
readelf -S throw_app | rg 'eh_frame|gcc_except_table'
readelf --debug-dump=frames throw_app | head -120
objdump -drwC -Mintel throw_app | rg 'landing|personality|__cxa_throw'
\end{lstlisting}

很多工程会规定热路径不抛异常，但“不在热路径抛异常”和“不理解异常机制”是两回事。即使你用错误码做业务错误，RAII 析构、栈回溯、崩溃报告、手写汇编边界和第三方库仍然会遇到 unwind。第一册必须把异常视为运行时合同，而不是语法糖。

\subsection{六、系统调用：libc wrapper、syscall、vDSO 和 page fault 边界}

系统调用不是普通函数调用。调用 \code{write} 时，源码层看起来是一个 C 函数；libc wrapper 内部可能处理 \code{errno}、线程取消点、重试和 ABI 适配；真正进入内核时，x86-64 Linux 使用 \instruction{syscall} 指令和 syscall ABI。普通函数 ABI 的第四个整数参数在 \register{rcx}，Linux syscall 常用 \register{r10}，并且 \instruction{syscall} 会破坏 \register{rcx} 和 \register{r11}。

\begin{center}
\small
\begin{tabular}{p{0.22\linewidth}p{0.34\linewidth}p{0.28\linewidth}}
\toprule
边界 & 发生了什么 & 证据 \\
\midrule
libc wrapper & 普通函数调用，处理库级语义 & 反汇编、符号、\cmd{ltrace} 或源码阅读 \\
syscall entry & 用户态切到内核态，使用 syscall ABI & \cmd{strace}、反汇编中的 \instruction{syscall} \\
vDSO & 某些查询走用户态映射页，避免真正陷入内核 & \filepath{/proc/self/maps}、\code{linux-vdso.so.1} \\
page fault & 指令访问虚拟地址时触发异常，内核修复或发信号 & \cmd{perf stat -e page-faults}、\cmd{strace -c}、maps \\
signal & 内核异步改变用户态控制流 & GDB signal、siginfo、core dump \\
\bottomrule
\end{tabular}
\end{center}

\code{clock\_gettime} 是理解 vDSO 的好例子。它有时不需要真正 syscall，而是通过内核映射到用户空间的 vDSO 页面读取时间相关数据。于是同样叫“系统时间”，成本可能远低于普通陷入内核。报告里若测量几十纳秒级 kernel，必须知道计时器本身走什么路径。

\begin{lstlisting}[numbers=none,caption={系统边界证据包}]
strace -e trace=write,read,mmap,brk,clock_gettime ./app
strace -c ./app
cat /proc/self/maps | rg 'vdso|vvar'
objdump -drwC -Mintel /lib/x86_64-linux-gnu/libc.so.6 | rg 'syscall'
perf stat -e page-faults,minor-faults,major-faults ./app
\end{lstlisting}

page fault 也要纳入系统调用边界的思维。它不是程序主动调用内核，而是 CPU 在取指、load 或 store 时发现页表状态不满足，转入异常处理。合法的首次触页、COW、文件映射缺页和非法地址崩溃，在机器层都从异常入口开始，但业务含义完全不同。若第一册不把 syscall、vDSO、page fault、signal 分清，后面做 benchmark 和服务 p99 分析一定会混层。

\subsection{七、signal：内核怎样异步改变用户态控制流}

信号是第一册必须单独建模的系统边界。它不像普通函数调用由当前源码位置主动发起，也不像 \code{syscall} 由用户态显式进入内核。信号可以来自硬件异常、终端、另一个进程、定时器、管道错误或内核事件；内核在合适的返回用户态时机，把用户线程的上下文改造成“先执行 signal handler，再回到原来位置或终止进程”。因此，signal 是异步控制流：它打断的可能是热循环、阻塞系统调用、库函数、内存分配器或 C++ 对象更新中间状态。

\begin{lstlisting}[numbers=none,caption={signal 的压缩状态机}]
UserRunning:
  thread executes ordinary user code or waits in a syscall

SignalGenerated:
  hardware exception, kill, terminal, timer, SIGPIPE or kernel event

Pending:
  signal is recorded for the process or a particular thread

DeliveryPoint:
  kernel chooses a thread and prepares user-space delivery

HandlerRunning:
  handler runs on the user stack or alternate signal stack

ReturnOrTerminate:
  sigreturn restores context, syscall may return EINTR,
  or default action terminates, stops or ignores the process
\end{lstlisting}

这个状态机解释两个常见现象。第一，程序收到 \code{SIGSEGV} 时，信号不是“随机函数回调”，而是 CPU 访问地址失败，内核无法修复 page fault，于是按信号机制报告给进程。第二，阻塞 \code{read}、\code{write}、\code{waitpid} 或 \code{futex} 期间收到信号，系统调用可能提前返回 \code{-1,EINTR}，也可能在 \code{SA\_RESTART} 和具体接口允许时被内核或 libc 重启。写系统代码时，\code{EINTR} 不是噪声，而是控制流边界。

\begin{center}
\small
\begin{tabular}{p{0.20\linewidth}p{0.33\linewidth}p{0.33\linewidth}}
\toprule
信号 & 常见来源 & 工程含义 \\
\midrule
\code{SIGSEGV} & 非法地址、权限错误、无法处理的 page fault & 看 fault address、\register{rip}、maps，不要只猜空指针 \\
\code{SIGBUS} & 某些映射文件访问、对齐或设备映射错误 & \code{mmap} 文件被截断时尤其要检查 \\
\code{SIGPIPE} & 向关闭的 pipe/socket 写入 & 服务端写路径要决定忽略、处理或用 \code{MSG\_NOSIGNAL} \\
\code{SIGINT}/\code{SIGTERM} & 终端或进程管理器请求退出 & 只能触发有序 shutdown，不能在 handler 里做复杂清理 \\
\code{SIGALRM}/定时器信号 & 计时器到期 & 会打断阻塞调用，benchmark 中要写明是否启用 \\
\bottomrule
\end{tabular}
\end{center}

安装 handler 应优先使用 \code{sigaction}，而不是把 signal 当成普通 C++ callback。handler 运行时的限制非常强：只能调用 async-signal-safe 的函数；不能随便 \code{new}、\code{delete}、锁 mutex、写 \code{std::string}、调用日志库或抛 C++ 异常。最稳的模式是 handler 只设置 \code{std::sig\_atomic\_t} 标志，或向专门 fd 写一个字节唤醒事件循环；真正清理工作回到普通控制流执行。

\begin{lstlisting}[language=C++,caption={signal handler 只做最小动作}]
#include <atomic>
#include <csignal>

static volatile std::sig_atomic_t stop_requested = 0;

extern "C" void on_signal(int signo) {
    if (signo == SIGTERM || signo == SIGINT) {
        stop_requested = 1;
    }
}

int install_signal_handlers() {
    struct sigaction sa {};
    sa.sa_handler = on_signal;
    sigemptyset(&sa.sa_mask);
    sa.sa_flags = 0;
    if (sigaction(SIGTERM, &sa, nullptr) != 0) return -1;
    if (sigaction(SIGINT, &sa, nullptr) != 0) return -1;
    return 0;
}
\end{lstlisting}

这个例子故意不在 handler 中打印日志。原因不是保守风格，而是运行时事实：信号可能打断正在持有 libc 内部锁、malloc 锁或日志锁的线程；handler 再调用同一套库函数，可能死锁或破坏状态。若需要把信号接入事件循环，Linux 工程常用 self-pipe trick、\code{eventfd} 或 \code{signalfd}，把异步信号转换成普通 fd 事件，再由主循环按普通状态机处理。

\subsubsection{signal frame：handler 为什么能回到被打断的位置}

信号处理函数看起来像普通函数：它有入口地址，执行完会返回。但它不是由用户代码里的 \instruction{call} 调用的。典型 Linux/x86-64 路径是：CPU 因异常、中断或系统调用返回路径进入内核；内核发现该线程有可投递信号；内核在用户栈上构造一个 \term{signal frame}{信号帧}，保存被打断时的寄存器、signal mask、故障信息和恢复所需的上下文；然后把即将返回用户态的 \register{rip} 改成 handler 地址，把用户栈调整到这个新 frame。handler 执行完后，不是普通地 \instruction{ret} 回到被打断函数，而是通过 C 运行时提供的 trampoline 发起 \code{rt\_sigreturn}，请求内核从 signal frame 恢复原来的寄存器、mask 和 \register{rip}。

\begin{lstlisting}[numbers=none,caption={signal delivery 的用户栈视角}]
Before delivery:
  user rsp -> ordinary user stack
  user rip -> instruction that was running or will resume

Kernel delivery:
  push/build rt signal frame on user stack
  save interrupted registers, fp state, mask, siginfo, ucontext
  set user rip = handler address
  set user rsp = signal frame top

Handler return:
  handler returns to trampoline
  trampoline performs rt_sigreturn syscall
  kernel validates frame and restores interrupted context
\end{lstlisting}

这解释了三个容易误判的事实。

第一，signal frame 不是 C++ 函数栈帧，但它占用用户栈。若线程栈已经接近耗尽，普通 handler 可能连 signal frame 都放不下，栈溢出本身的 \code{SIGSEGV} handler 也可能再次崩溃。需要可靠处理栈溢出、崩溃记录或 profiling 信号时，应考虑 \code{sigaltstack} 和 \code{SA\_ONSTACK}，让 handler 在预先准备的备用栈上运行。

第二，handler 打断的位置可能是任意普通控制流边界。它可能出现在一条 C++ 语句中间、一个 \code{malloc} 内部、一个日志库持锁区、一个无锁结构的更新窗口，或一个阻塞系统调用返回前。源码层“我没有调用这个函数”的直觉没有意义；信号投递是内核改写返回用户态上下文，不是源代码显式调用。

第三，signal frame 是安全边界。内核恢复上下文前要验证 frame；调试器和崩溃分析也要区分普通调用栈帧、signal trampoline 和被打断帧。若 unwinder、profiler 或手写汇编不能识别 signal frame，就可能把栈回溯切断，误以为 handler 是从某个普通函数直接调用出来的。

\begin{center}
\small
\begin{tabular}{p{0.22\linewidth}p{0.34\linewidth}p{0.30\linewidth}}
\toprule
栈帧类型 & 由谁建立 & 调试重点 \\
\midrule
C++ 普通栈帧 & 编译器按 ABI 生成 prologue/epilogue & 返回地址、callee-saved 寄存器、局部变量位置 \\
系统调用内核栈 & CPU/内核在进入内核后使用 & 用户态通常看不到，只能通过 trace/perf 间接观察 \\
signal frame & 内核在用户栈或 alternate stack 上构造 & \code{siginfo}、\code{ucontext}、被打断 \register{rip}、\code{rt\_sigreturn} \\
\bottomrule
\end{tabular}
\end{center}

\subsubsection{mask、pending 和 delivery：信号不是普通队列}

信号还有一层状态经常被忽略：每个线程都有 signal mask，进程和线程有 pending 信号集合，某些信号的处理方式是进程级 disposition。一个信号生成后，若当前目标线程屏蔽了它，它不会立刻运行 handler，而是变成 pending；等 mask 解除或内核选择了可投递线程，才进入 delivery。传统非实时信号通常按“是否 pending”合并，多个同类信号不一定形成多个排队事件；实时信号才有更接近排队的语义。

\begin{lstlisting}[numbers=none,caption={signal mask 和 pending 的状态机}]
Unblocked:
  signal can be delivered at a return-to-user or interruptible boundary

Blocked:
  generated signal becomes pending, handler does not run yet

Pending:
  signal is remembered for a process or thread

Unmasked:
  kernel may select a target thread and build a signal frame

HandledOrDefaulted:
  handler runs, or default action terminates/stops/ignores
\end{lstlisting}

多线程程序里还要分清“发给进程”和“发给线程”。\code{kill(pid, signo)} 面向进程；\code{pthread\_kill(thread, signo)} 面向线程；硬件异常通常属于触发异常的那个线程。一个服务若有 worker 线程池，最好明确哪些线程屏蔽终止信号，哪个控制线程通过 \code{sigwait}、\code{signalfd} 或 event loop 统一接管。否则同一个 \code{SIGTERM} 可能落在任意未屏蔽线程上，handler 打断的位置也就不可预测。

\begin{lstlisting}[numbers=none,caption={观察 mask、pending 和投递路径}]
# /proc 中的十六进制位图能说明哪些信号被阻塞、忽略、捕获或 pending。
cat /proc/$PID/status | grep -E 'SigBlk|SigIgn|SigCgt|SigPnd|ShdPnd'

# strace 能看到安装 handler、改变 mask 和信号实际投递。
strace -f -ttT -e trace=rt_sigaction,rt_sigprocmask,rt_sigsuspend,rt_sigreturn ./app
\end{lstlisting}

这也是为什么“临时屏蔽信号保护临界区”不是通用同步方案。它只能影响信号投递，不能阻止其他线程访问共享数据，也不能替代 mutex、atomic 或内存序。信号 mask 是异步控制流工具，不是普通并发互斥工具。

\subsubsection{async-signal-safe：根因是可重入性，不是函数表迷信}

POSIX 列出 async-signal-safe 函数表，但背后的根因更重要：handler 可以在任意指令边界重入当前线程。如果普通控制流正在执行一个非可重入函数，内部状态可能处于“锁已持有、链表改到一半、buffer 指针暂时不一致、TLS 缓存正在更新”的中间态；handler 再调用同一个库或依赖同一把锁，就可能死锁或破坏状态。

\begin{center}
\small
\begin{tabular}{p{0.25\linewidth}p{0.32\linewidth}p{0.29\linewidth}}
\toprule
handler 中的动作 & 风险 & 更稳的替代 \\
\midrule
\code{printf}/日志库 & stdio 或日志锁可能已被打断线程持有 & \code{write} 固定 fd，或只设置标志 \\
\code{malloc}/\code{new} & allocator 内部状态可能正在更新 & 预分配 buffer，或回到普通控制流处理 \\
\code{pthread\_mutex\_lock} & handler 可能重入同一线程已持有的锁 & 使用 \code{sig\_atomic\_t} 标志或 self-pipe \\
抛 C++ 异常 & 跨 signal frame 展开不是普通 C++ 调用语义 & 记录请求，普通控制流检查后退出 \\
修改复杂对象 & 可能打断对象不变量维护窗口 & 只写原子标志或单字节通知 \\
\bottomrule
\end{tabular}
\end{center}

即使使用 \code{std::atomic}，也不能把 handler 当成普通线程函数。C++ 标准、POSIX、libc 实现和硬件原子之间有边界；最可移植的最小模式仍是 \code{volatile sig\_atomic\_t} 标志，或 Linux 工程中常见的 \code{write} 到 pipe/eventfd，让主循环在普通上下文里完成复杂工作。追求“handler 里直接做完清理”通常会把 shutdown 写成偶发死锁源。

\subsubsection{EINTR 和 SA\_RESTART：阻塞调用怎样被信号切开}

信号也会把“普通阻塞控制流”切开。线程可能正在 \code{read}、\code{ppoll}、\code{accept}、\code{waitpid}、\code{futex} 或带超时的等待里睡眠；信号到达后，内核需要先让 handler 运行。handler 返回后，原系统调用有三类结果：已经完成并返回成功；返回 \code{-1} 且 \code{errno=EINTR}；或在 \code{SA\_RESTART} 和接口语义允许时自动重启。

\begin{lstlisting}[numbers=none,caption={被信号打断的阻塞调用不能机械重试}]
BlockingSyscall:
  thread sleeps in read, ppoll, waitpid, futex, accept, ...

SignalDelivered:
  handler runs before syscall result is reported to user code

AfterHandler:
  syscall may have completed
  syscall may fail with EINTR
  syscall may be restarted if SA_RESTART and the API allow it

CallerPolicy:
  retry, recompute timeout, handle cancellation, or report partial progress
\end{lstlisting}

因此，处理 \code{EINTR} 的正确问题不是“要不要 while 重试”，而是“这个接口的业务语义是什么”。读写类调用可能已经发生部分进度；带 deadline 的等待必须重算剩余时间；取消语义可能正是通过信号触发；\code{ppoll}/\code{pselect} 还可能故意把 mask 切换和等待合成一个原子边界，避免“解除屏蔽后、真正睡眠前”丢信号。服务端 p99 报告若出现信号、超时和重试，必须把这些路径写进 trace，而不是只看用户态热循环。

\subsubsection{备用信号栈：为什么栈溢出 handler 不能依赖原栈}

默认情况下，handler 使用被打断线程的用户栈。大多数终止信号处理和轻量通知可以接受这个事实；但栈溢出、深递归、坏指针破坏栈、profiling 信号打断极深调用链时，原栈可能已经不可靠。 \code{sigaltstack} 允许线程注册一块独立内存，配合 \code{SA\_ONSTACK} 让指定 handler 在备用栈上运行。

\begin{lstlisting}[language=C++,caption={备用 signal stack 的最小骨架}]
#include <signal.h>
#include <stdlib.h>

static stack_t alt_stack;

int install_alt_stack() {
    alt_stack.ss_sp = malloc(SIGSTKSZ);
    if (alt_stack.ss_sp == nullptr) return -1;
    alt_stack.ss_size = SIGSTKSZ;
    alt_stack.ss_flags = 0;
    return sigaltstack(&alt_stack, nullptr);
}

int install_crash_handler(void (*handler)(int, siginfo_t*, void*)) {
    struct sigaction sa {};
    sa.sa_sigaction = handler;
    sigemptyset(&sa.sa_mask);
    sa.sa_flags = SA_SIGINFO | SA_ONSTACK;
    return sigaction(SIGSEGV, &sa, nullptr);
}
\end{lstlisting}

这个骨架只说明机制，不代表可以在崩溃 handler 里安全做复杂诊断。崩溃时 heap、锁、TLS、stdio 和 C++ 对象都可能处于损坏状态；高质量 crash reporter 通常预分配资源，用 async-signal-safe 的最小写路径记录寄存器、fault address、线程 id 和少量上下文，然后让进程终止或交给外部收集 core dump。

\begin{lstlisting}[numbers=none,caption={signal 证据包}]
# 观察崩溃信号和 fault address。
gdb -q ./app core
(gdb) info signals
(gdb) p $_siginfo
(gdb) info registers
(gdb) x/i $rip
(gdb) bt

# 观察系统调用是否被信号打断。
strace -f -ttT -e trace=read,write,ppoll,rt_sigaction,rt_sigprocmask,rt_sigreturn ./app

# 观察当前进程的 signal mask 和 pending 状态。
cat /proc/$PID/status | grep -E 'SigBlk|SigIgn|SigCgt|SigPnd|ShdPnd'

# 观察 core dump 中 signal trampoline 和被打断帧。
eu-stack -p $PID
coredumpctl gdb ./app
\end{lstlisting}

signal 报告也要写证据边界。看到 \code{SIGSEGV}，只能说明进程因为内存访问或相关异常被终止，不能自动推出“空指针”。要结合 fault address、\register{rip} 指令、寄存器、\filepath{/proc/<pid>/maps}、对象生命周期和 sanitizer 结果。看到 \code{EINTR}，只能说明某个阻塞边界被信号打断，不能自动重试所有接口；有些调用已经部分完成，有些调用有超时语义，有些业务正是用信号请求取消。看到 \code{SIGPIPE}，也不能只把它关掉：要先说明写入对象是 pipe、socket 还是普通文件，对端关闭时业务应该丢弃、重连、还是报告失败。

\subsection{深水链路的统一验收}

本节不是要求读者立刻成为 LLVM、ELF、DWARF 或内核专家，而是要求形成工程闭环意识。一个合格的第一册实验报告，至少要能回答：

\begin{itemize}
  \item 这段源码在 AST/IR 层的核心控制流和数据流是什么。
  \item 哪些 pass 改写了它，哪些 pass 被别名、异常、副作用或目标 ISA 阻止。
  \item 最终汇编是否和 IR 结论一致，是否出现意外 PLT、间接调用或栈槽。
  \item ELF 中是否存在 \code{.init\_array}、TLS、动态段和需要运行时绑定的符号。
  \item 调试器看到的变量、栈和源码行是否由 DWARF/CFI 支持。
  \item 异常路径是否有 unwind table、personality、cleanup 和 RAII 证据。
  \item 计时范围是否混入 syscall、vDSO、page fault、signal、动态链接或初始化成本。
\end{itemize}

能回答这些问题，第一册才真正从“看懂汇编”进入“解释 C++ 程序如何被机器、运行时和操作系统共同执行”的层次。

\section{反例：看似合理但错误的理解}

\begin{warningbox}
\textbf{误区一：C++ 程序从 \code{main} 的第一行开始。}

不准确。进程入口通常是 runtime 启动代码，\code{main} 是被运行时调用的用户入口。

\textbf{误区二：头文件会被单独编译。}

通常不对。头文件被包含进源文件，形成翻译单元。

\textbf{误区三：目标文件就是可执行文件。}

不对。目标文件还需要链接。

\textbf{误区四：源码里有函数，机器码里一定有这个函数。}

不对。优化后函数可能被 inline、删除、合并或改变可见性。

\textbf{误区五：指针就是物理地址。}

用户态程序看到的是虚拟地址。
\end{warningbox}

\section{作业}

\begin{exercisebox}
\textbf{作业 1：画出完整流程。}

画出从 \code{hello.cpp} 到正在运行进程的完整路径。每一层写明：

\begin{itemize}
  \item 输入是什么。
  \item 输出是什么。
  \item 哪个工具或系统组件负责。
  \item 性能学习为什么要关心。
\end{itemize}

\textbf{作业 2：函数消失实验。}

写三个函数：

\begin{itemize}
  \item 一个允许 inline。
  \item 一个使用 \code{noinline} 阻止 inline。
  \item 一个从不调用。
\end{itemize}

分别用 \code{-O0} 和 \code{-O3} 编译，用 \code{nm} 和 \code{objdump} 观察哪些符号存在，哪些函数消失。

\textbf{作业 3：启动成本和 kernel 成本。}

写一个程序，让 \code{main} 调用某个小函数一次，再调用同一函数一亿次。分别测整个程序时间和循环区域时间。解释为什么启动成本不能代表 kernel 成本。
\end{exercisebox}

\section{验收标准}

本章收束时，应能：

\begin{itemize}
  \item 解释预处理、编译、汇编、链接、加载的区别。
  \item 生成并解释 \code{.i}、\code{.s}、\code{.o} 和可执行文件。
  \item 用 \code{readelf}、\code{nm}、\code{objdump} 做基础观察。
  \item 解释为什么 \code{main} 不是真正的进程入口。
  \item 画出简化虚拟地址空间。
  \item 说明为什么性能测量要排除启动、初始化、动态链接和全局构造等因素。
\end{itemize}
